\documentclass[12pt,oneside]{uhthesis}
\usepackage{subfigure}
\usepackage[ruled,lined,linesnumbered,titlenumbered,algochapter,spanish,onelanguage]{algorithm2e}
\usepackage{amsmath}
\usepackage{amssymb}
\usepackage{amsbsy}
\usepackage{caption,booktabs,array}
\captionsetup{ justification = centering }
%\usepackage{mathpazo}
\usepackage{float}
\setlength{\marginparwidth}{2cm}
\usepackage{todonotes}
\usepackage{listings}
\usepackage{xcolor}
\usepackage{multicol}
\usepackage{graphicx}
% Diagramas de la tesis (Graphics/*.tex).
\usepackage{tikz}
\usetikzlibrary{arrows.meta,backgrounds,fit}
% Permite partir URL largas (bibliografía).
\usepackage{xurl}
\floatstyle{plaintop}
\restylefloat{table}
\addbibresource{Bibliography.bib}
% Figuras de la tesis y las que regenera P09 en redaccion/figuras/.
\graphicspath{{Graphics/}{../figuras/}}
% \setlength{\parskip}{\baselineskip}%
\renewcommand{\tablename}{Tabla}
\renewcommand{\listalgorithmcfname}{Índice de Algoritmos}
%\dontprintsemicolon
\SetAlgoNoEnd

\definecolor{codegreen}{rgb}{0,0.6,0}
\definecolor{codegray}{rgb}{0.5,0.5,0.5}
\definecolor{codepurple}{rgb}{0.58,0,0.82}
\definecolor{backcolour}{rgb}{0.95,0.95,0.92}

\lstdefinestyle{mystyle}{
    backgroundcolor=\color{backcolour},   
    commentstyle=\color{codegreen},
    keywordstyle=\color{purple},
    numberstyle=\tiny\color{codegray},
    stringstyle=\color{codepurple},
    basicstyle=\ttfamily\footnotesize,
    breakatwhitespace=false,         
    breaklines=true,                 
    captionpos=b,                    
    keepspaces=true,                 
    numbers=left,                    
    numbersep=5pt,                  
    showspaces=false,                
    showstringspaces=false,
    showtabs=false,                  
    tabsize=4
}

\lstset{style=mystyle}

% Título provisional (R00, 2026-11-01): traducción del nombre del repositorio.
\title{Generación narrativa bidireccional: planificación \textit{Top-Down} y emergencia \textit{Bottom-Up}}
\author{Josue Rolando Naranjo Sieiro}
\advisor{Dr. Alejandro Piad Morffis}
\advisorlabel{Tutor}
\degree{Licenciado en Ciencia de la Computación}
\faculty{Facultad de Matemática y Computación}
% El enlace del repositorio, partido a mano en dos líneas que apuntan a la misma URL.
\newcommand{\repourl}{https://github.com/JD3M0N/BiPlan-Forge-Bidirectional-Narrative-Generation-through-Top-Down-Planning-and-Bottom-Up-Emergence}
\date{Noviembre de 2026\\{\urlstyle{same}%
  \href{\repourl}{\nolinkurl{github.com/JD3M0N/BiPlan-Forge-Bidirectional-Narrative-Generation-}}\\
  \href{\repourl}{\nolinkurl{through-Top-Down-Planning-and-Bottom-Up-Emergence}}}}
\logo{Graphics/uhlogo}
\makenomenclature

\renewcommand{\vec}[1]{\boldsymbol{#1}}
\newcommand{\diff}[1]{\ensuremath{\mathrm{d}#1}}
\newcommand{\me}[1]{\mathrm{e}^{#1}}
\newcommand{\pf}{\mathfrak{p}}
\newcommand{\qf}{\mathfrak{q}}
%\newcommand{\kf}{\mathfrak{k}}
\newcommand{\kt}{\mathtt{k}}
\newcommand{\mf}{\mathfrak{m}}
\newcommand{\hf}{\mathfrak{h}}
\newcommand{\fac}{\mathrm{fac}}
\newcommand{\maxx}[1]{\max\left\{ #1 \right\} }
\newcommand{\minn}[1]{\min\left\{ #1 \right\} }
\newcommand{\lldpcf}{1.25}
\newcommand{\nnorm}[1]{\left\lvert #1 \right\rvert }
\renewcommand{\lstlistingname}{Ejemplo de código}
\renewcommand{\lstlistlistingname}{Ejemplos de código}

\begin{document}

\frontmatter
\maketitle

% Dedicatoria, agradecimientos y opinión del tutor: todavía no (ver ../write_thesis_roadmap.md).
% \begin{dedication}
    Dedicación
\end{dedication}
% \begin{acknowledgements}
    Agradecimientos
\end{acknowledgements}
% \begin{opinion}
    Opiniones de los tutores
\end{opinion}
\begin{resumen}
% R10 en ../write_thesis_roadmap.md
\end{resumen}

\begin{abstract}
% R10 en ../write_thesis_roadmap.md
\end{abstract}

\tableofcontents
\listoffigures
\listoftables
% \listofalgorithms
% \lstlistoflistings


\mainmatter

\begin{introduction}\label{chapter:introduction}
% R02 en ../write_thesis_roadmap.md. Sin referencias nuevas: todas las claves están en
% Bibliography.bib desde R01. Hipótesis H1-H3 de docs/resumen.md; la forma de H2 y H3 la
% decidió Josue el 2026-11-01 (H2 se mantiene sin experimento; H3 como lectura de WSE-bench).

Generar historias de forma automática es uno de los problemas más antiguos de la inteligencia
artificial. TALE-SPIN~\cite{meehan1977}, a finales de los años setenta, ya producía fábulas
razonando cómo cada personaje resolvía sus objetivos. Durante décadas, los generadores dependieron
de mundos descritos a mano en lenguajes formales, que garantizaban la causalidad de los eventos a
costa de acotar las historias posibles~\cite{riedl2010}. Los \textit{Large Language Models}
(LLMs), modelos de lenguaje de gran tamaño capaces de seguir instrucciones en lenguaje natural,
cambiaron ese límite: producen prosa fluida a partir de una petición, sin un modelo del dominio
programado de antemano. Lo que no garantizan es que esa prosa sea una buena historia, y tampoco
ofrecen una forma fiable de saber si lo es. Esta tesis trata de esos dos problemas: generar
historias con LLMs y evaluarlas.

La literatura ordena los generadores según quién decide la trama. En la generación dirigida por la
trama (\textit{Top-Down}), un componente con visión de autor decide los eventos y los personajes
los ejecutan; con LLMs, eso se traduce en planificar la historia antes de
escribirla~\cite{yang2023,huot2025}. En la generación por simulación de personajes
(\textit{Bottom-Up}), la trama no existe de antemano: emerge de lo que hacen unos agentes con
objetivos y memoria propios~\cite{aylett1999,park2023}. Entre ambas hay una tradición de sistemas
híbridos, como Façade~\cite{mateas2005} o The Virtual Storyteller~\cite{theune2003}, que combinan
un componente de autor con personajes autónomos y que los LLMs han reabierto~\cite{chen2026a}. El
Capítulo~\ref{chapter:state-of-the-art} precisa estos términos y recorre las tres familias.

El antecedente directo de este trabajo es la tesis de Fuentes Rodríguez~\cite{fuentes2025},
presentada en esta facultad en 2025 y cotutorada por el mismo tutor. Su sistema generaba
historias mediante la simulación de agentes conversacionales y lo comparaba con un modelo guiado
por un solo \textit{prompt} (la instrucción en lenguaje natural que recibe el modelo). Doce
voluntarios puntuaron ambas salidas de 1 a 10, y los perfiles resultantes apenas se distinguieron.
De ese trabajo se conservan la línea base de un solo \textit{prompt} y el interés por la coherencia
de las historias largas; se cambia la forma de evaluar.

Los problemas que motivan esta tesis se agrupan en tres:

\begin{itemize}
  \item \textbf{La paradoja narrativa:} la simulación de personajes da comportamientos creíbles,
  pero nadie vela por la estructura que convierte una serie de eventos en una historia; la
  planificación vela por esa estructura, pero tiende a reducir a los personajes a ejecutores del
  plan~\cite{louchart2003,riedl2010}. Un generador que use solo una de las dos renuncia a lo que
  aporta la otra.
  \item \textbf{La pérdida de compromisos y de riqueza:} aun con el contexto necesario delante, los
  LLMs dejan sin pagar muchas de las prefiguraciones que introducen~\cite{yun2026}, se contradicen a
  lo largo de una historia larga~\cite{li2026} y hacen que un personaje revele lo que en ese momento
  no podía saber~\cite{ahn2024}. Razonar quién sabe qué en una conversación les resulta mucho más
  difícil que a las personas~\cite{kim2023}, y sus historias superan menos pruebas de creatividad
  que las de escritores profesionales~\cite{chakrabarty2024}.
  \item \textbf{La evaluación poco fiable:} las escalas de valoración absolutas son las más usadas y
  las más criticadas~\cite{amidei2019}, y con ellas incluso evaluadores expertos coinciden
  poco~\cite{chiang2023}. Los LLMs usados como jueces favorecen sus propias generaciones y las
  historias escritas por IA~\cite{zheng2023,panickssery2024,tutone2026}. Sin una evaluación fiable
  no se puede saber si un mecanismo nuevo mejora la historia o solo la cambia.
\end{itemize}

El tema es actual. Los sistemas híbridos con LLMs son de los últimos años: StoryBox~\cite{chen2026a}
y WSE-bench~\cite{chen2026b}, este último todavía un preprint, combinan agentes simulados con un
narrador, y Yu et al.~\cite{yu2025}, con un plan. También son recientes el trabajo de Yun et al.~\cite{yun2026} sobre los
compromisos sin pagar y los resultados que cuestionan a los jueces LLM en escritura creativa o que
proponen aprenderlos de preferencias humanas~\cite{tutone2026,fein2026}. Además, las historias de
esta tesis están en español, una lengua en la que, en una comparación de GPT-4 con un novelista, la
brecha entre ambos fue mayor que en inglés~\cite{marco2024}.

La \textbf{novedad científica} del trabajo está en dos aportaciones. La primera es un generador
híbrido, bidireccional en el sentido del título: el plan baja de arriba abajo
(\textit{Top-Down}) y los personajes simulados suben de abajo arriba (\textit{Bottom-Up}). Separa tres
decisiones. Un plan validado decide qué pasa; unos personajes simulados que
solo saben lo que presenciaron deciden cómo pasa, al representar el plan escena a escena en lo que
llamamos una \textbf{función}, como en el teatro; y un narrador que cura el registro de esa función,
sin inventar, decide cómo se cuenta. En la literatura que revisa el
Capítulo~\ref{chapter:state-of-the-art}, que no es exhaustiva, no encontramos un generador con LLMs
que reúna las tres piezas tal como aquí se formulan; los dos más cercanos, Yu et al. y PLOTTER,
comparten solo parte de ellas~\cite{yu2025,gu2026} (Sección~\ref{sec:posicion}). La segunda es una
metodología de evaluación en la que el LLM no juzga la calidad: cuenta rasgos de la historia y cita el
pasaje de cada recuento, que el código verifica literalmente. La calidad la deciden lectores que
comparan pares de historias a ciegas, y un juez aprendido de esas comparaciones intenta ordenar las
historias como ellos a partir de los recuentos, con una validación fuera de muestra que dice si lo
consigue.

La importancia teórica del trabajo es doble. Como el plan, la función y la narración son etapas
separadas, el generador sirve de banco de pruebas de la paradoja narrativa: se puede cambiar una
etapa y conservar las demás, por ejemplo escribir la historia directamente desde el mismo plan o
narrar la misma función desde otro punto de vista, y medir qué aporta cada una. La metodología, por
su parte, muestra cómo evaluar historias sin pedir a un modelo un juicio de calidad, una cuestión
abierta en el área (Sección~\ref{sec:evaluacion-historias}). En la práctica, el generador y las
herramientas de evaluación forman un sistema que recalcula sus resultados desde los datos guardados,
que residen en la copia del autor y no en Git (Sección~\ref{sec:pruebas}), con un bot de Telegram para pedir
historias y una interfaz local para leerlas a ciegas, y el juez lee solo el texto de la historia,
de modo que la metodología se puede aplicar a las de otro generador.

A partir de lo anterior se plantea el siguiente \textbf{problema}: ¿cómo generar historias con LLMs
que conserven los compromisos de su trama sin reducir a los personajes a ejecutores del plan, y cómo
evaluarlas de forma fiable sin pedir a un modelo que juzgue su calidad? El \textbf{objeto de
estudio} es la generación automática de historias con modelos de lenguaje, y el \textbf{campo de
acción}, la generación híbrida de historias en español, que combina la planificación de la trama
con la simulación de personajes, junto con su evaluación mediante rasgos contables y comparaciones
humanas por pares.

Para resolver el problema se plantea como \textbf{Objetivo General} desarrollar un generador
híbrido de historias con LLMs que combine un plan validado con la simulación de personajes con
memoria propia, junto con una metodología para evaluarlo basada en recuentos verificables y
comparaciones humanas por pares a ciegas. Para lograrlo se plantean los siguientes
\textbf{Objetivos Específicos}:

% Correspondencia objetivo -> capítulo: 1 -> Cap. 1 (Background.tex), 2 -> Cap. 2
% (Generator.tex), 3 -> Cap. 3 (Evaluation.tex), 4 -> Cap. 4 (Implementation.tex),
% 5 -> Cap. 5 (Experiments.tex). R08 escribe una conclusión por objetivo, en este orden.
\begin{itemize}
  \item \textbf{Analizar el estado del arte} de la generación de historias dirigida por la trama,
  por simulación de personajes e híbrida, y de su evaluación, para identificar qué adoptar y qué
  hueco queda. % -> Capítulo 1
  \item \textbf{Diseñar un generador híbrido} en el que un plan validado decida qué pasa, una
  función con personajes de memoria propia decida cómo pasa y un narrador decida cómo se
  cuenta. % -> Capítulo 2
  \item \textbf{Diseñar una metodología de evaluación} que combine rasgos contables con cita
  verificada, comparaciones humanas por pares a ciegas que incluyan una línea base de un solo
  \textit{prompt}, y un juez aprendido de esas comparaciones y validado fuera de
  muestra. % -> Capítulo 3
  \item \textbf{Implementar el generador y las herramientas de evaluación} en un sistema
  reproducible, con una interfaz para la lectura a ciegas. % -> Capítulo 4
  \item \textbf{Evaluar experimentalmente el generador}: validarlo antes de fijar su configuración,
  validar la extracción de rasgos, comparar parametrizaciones, recoger las preferencias de los
  lectores y ordenar el corpus de historias. % -> Capítulo 5
\end{itemize}

Las \textbf{hipótesis} de trabajo son tres. Las dos primeras comparan la configuración del
generador con una alternativa sobre la misma entrada; la tercera, cada mecanismo opcional con su
ausencia. Todas se contrastan con recuentos de rasgos y con la preferencia de los lectores en
comparaciones a ciegas, según los tres criterios de la evaluación: creatividad, desarrollo de los
personajes e interés de la trama.

% Correspondencia hipótesis -> experimentos (R06 y R07 en ../write_thesis_roadmap.md):
% H1 -> R06: celdas M03 (narrativo --plan-from) frente a su base M01 en la matriz A (EXP-1),
%       parcial: 14 de 35 unidades (P05 cerrada el 2026-10-30). R07: pares narrativo-simulado
%       del estudio humano (solo el piloto; P06 cerrada el 2026-10-31 sin estudio real).
% H2 -> R07: EXP-2 (--actor-memory own frente a shared) no se corrió (idea-bank desde P05.1,
%       2026-10-30); respuesta prevista: sin datos suficientes. Evidencia indirecta en R06: la
%       auditoría de 62 contextos de actor de la validación 7.5 (aislamiento, no mejora).
% H3 -> R06: validaciones antes de congelar el generador, P02.1 (narración, empate 2-2-1 en 5
%       pares) y P02.2 (brazos Top-Down, A 3 - B 1 en 4 pares), con un solo lector; R07, si el
%       juez o el piloto añaden algo.
% R12 (2026-11-03): las tres se responden en Experiments.tex, sec:respuesta-hipotesis, «sin datos
% suficientes» (en H3, lo observado es coherente con lo que anticipa: P02.1, P02.2 y 24-24 frente a la línea base en el piloto).
\begin{itemize}
  \item \textbf{H1:} con un mismo plan, la historia narrada a partir de una función difiere de la
  escrita directamente desde el plan en rasgos contables de personajes y de trama, más allá del
  ruido entre repeticiones, y los lectores la prefieren a ciegas en al menos uno de los tres
  criterios.
  \item \textbf{H2:} dar a cada personaje solo la memoria de lo que presenció, en lugar de la de
  todo lo ocurrido en público, produce escenas con menos menciones a hechos que el personaje no
  podía conocer, y que los lectores prefieren a ciegas.
  \item \textbf{H3:} más mecanismos no implican una historia mejor: que un mecanismo mejore medidas
  de la prosa no basta para que los lectores prefieran a ciegas las historias que lo usan frente a
  las que no. Traslada parte de la conclusión de WSE-bench~\cite{chen2026b}, que mide la coherencia
  con jueces LLM, a la preferencia de los lectores en este generador, y obliga a que cada mecanismo se
  gane su complejidad.
\end{itemize}

Las hipótesis no fijan de antemano qué resultado las confirmaría o refutaría. La regla de ruido de la
Sección~\ref{sec:horizontal} decide solo si una diferencia de la matriz se interpreta para un rasgo,
no cuántos rasgos tendrían que superarla para dar por cumplida H1, y la parte humana de las
hipótesis no tiene umbral.

El documento se organiza en cinco capítulos. El Capítulo~\ref{chapter:state-of-the-art} revisa el
estado del arte de la generación y la evaluación de historias y fija la posición del trabajo. El
Capítulo~\ref{chapter:generator} describe el generador híbrido: el plan, la función, el narrador y
los mecanismos opcionales. El Capítulo~\ref{chapter:evaluation} expone la metodología de
evaluación, y el Capítulo~\ref{chapter:implementation}, la implementación del sistema y sus
herramientas. El Capítulo~\ref{chapter:experiments} presenta los experimentos, con su alcance y sus
resultados incompletos, y responde a las hipótesis con lo que los datos permiten. Por último, las
Conclusiones recogen lo alcanzado en cada objetivo, las aportaciones y las limitaciones, y las
Recomendaciones, las líneas de trabajo futuro. Los anexos dan la lista núcleo de rasgos, las
peticiones de las obras de los experimentos y dos tablas complementarias.
\end{introduction}

\chapter{Estado del arte}\label{chapter:state-of-the-art}
% R01 en ../write_thesis_roadmap.md

\chapter{Generador híbrido de historias}\label{chapter:generator}
% R03 en ../write_thesis_roadmap.md. Fuentes: docs/resumen.md, docs/architecture/generation.md y
% simulation.md, y el código de referencia en Stagecraft 7.22 (etiqueta generator-frozen,
% 2026-10-28; predeterminados en planning/decisions.md §3 y options.py). Referencia nueva:
% piad2026 (B02 de la planilla).

Este capítulo describe el generador de historias propuesto, con la configuración de referencia
establecida para los experimentos el 28 de octubre de 2026. Esta referencia permite identificar las
condiciones de las ejecuciones mientras continúa el desarrollo. Primero se presenta la
\textbf{Arquitectura} (Sección~\ref{sec:arquitectura}), que separa en etapas las tres decisiones de
una historia: qué pasa, cómo pasa y cómo se cuenta. Después se describe cada etapa: el \textbf{Plan
validado} y su contrato de promesas (Sección~\ref{sec:plan}), el \textbf{Guion teatral}
(Sección~\ref{sec:guion}), la \textbf{Función simulada}, con el reparto, la memoria de los actores
y el director (Sección~\ref{sec:funcion}), y el \textbf{Narrador} (Sección~\ref{sec:narrador}). La
Sección~\ref{sec:opcionales} reúne los \textbf{Mecanismos opcionales}, que el generador implementa
pero que su configuración de referencia deja apagados.

\section{Arquitectura}\label{sec:arquitectura}

Como se muestra en la Figura~\ref{fig:arquitectura}, el generador encadena etapas, cada una a cargo
de un agente, es decir, de una llamada a un LLM con su propio \textit{prompt} y su propio formato de
salida. Las primeras etapas deciden qué pasa desde la perspectiva del autor: a partir de la
petición se eligen la escala de la historia, el mundo y los personajes; se construye un plan que el
código valida como un grafo de eventos, se traza sobre él un contrato de promesas y un dramaturgo lo
convierte en un guion de actos y escenas. Es la mitad \textit{Top-Down}. En la mitad
\textit{Bottom-Up}, unos actores simulados representan cada escena sin conocer el guion, y de su
interacción sale cómo pasa: es la función, cuyo único producto es el registro de lo que cada
personaje dijo, hizo y pensó. Por último, un narrador convierte ese registro en la historia y
decide cómo se cuenta, es decir, el discurso en el sentido de Chatman
(Capítulo~\ref{chapter:state-of-the-art}).

\begin{figure}[t]
  \centering
  % Arquitectura del generador (R03, Figura fig:arquitectura). Se incluye con \input desde
% MainMatter/Generator.tex; necesita tikz con arrows.meta, backgrounds y fit (Thesis.tex).
\begin{tikzpicture}[
    font=\footnotesize,
    caja/.style={draw, rounded corners=2pt, align=center, text width=2.2cm,
      minimum height=1.25cm, inner sep=2pt, fill=white},
    flecha/.style={-{Stealth[length=1.8mm]}, semithick},
    banda/.style={draw=black!45, dashed, rounded corners=4pt, inner sep=4pt},
    rotulo/.style={font=\footnotesize\itshape, text=black!70},
  ]
  % Primera fila: qué pasa.
  \node[caja] (peticion) at (0,0) {Petición};
  \node[caja] (analisis) at (2.9,0) {Análisis: perfil, mundo y personajes};
  \node[caja] (plan) at (5.8,0) {Plan: grafo de eventos validado};
  \node[caja] (promesas) at (8.7,0) {Contrato de promesas};
  \node[caja] (guion) at (11.6,0) {Guion: actos y escenas};
  % Segunda fila: cómo pasa y cómo se cuenta.
  \node[caja] (reparto) at (0,-2.75) {Reparto y compuertas};
  \node[caja] (funcion) at (2.9,-2.75) {Función: director y actores};
  \node[caja] (registro) at (5.8,-2.75) {Registro de la función};
  \node[caja] (narrador) at (8.7,-2.75) {Narrador};
  \node[caja] (historia) at (11.6,-2.75) {Historia};

  \foreach \a/\b in {peticion/analisis, analisis/plan, plan/promesas, promesas/guion,
      reparto/funcion, funcion/registro, registro/narrador, narrador/historia}
    \draw[flecha] (\a) -- (\b);
  \draw[flecha] (guion.south) -- ++(0,-0.75) -| (reparto.north);
  % Formato narrativo: la historia se escribe directamente desde el plan y las promesas.
  \draw[flecha, dashed] (promesas.north) -- ++(0,0.55) -| ([xshift=4.5mm]historia.east)
    node[rotulo, pos=0.25, above] {formato narrativo}
    -- (historia.east);

  \begin{scope}[on background layer]
    \node[banda, fit=(peticion)(guion),
      label={[rotulo, anchor=south west]north west:\textit{Top-Down}: qué pasa}] {};
    \node[banda, fit=(reparto)(registro),
      label={[rotulo, anchor=north west]south west:\textit{Bottom-Up}: cómo pasa}] {};
    \node[banda, fit=(narrador)(historia),
      label={[rotulo, anchor=north west]south west:Discurso: cómo se cuenta}] {};
  \end{scope}
\end{tikzpicture}

  \caption{Arquitectura del generador: qué pasa, cómo pasa y cómo se cuenta}
  \label{fig:arquitectura}
\end{figure}

La misma cadena admite tres formatos de salida, que comparten todo lo anterior al punto en que se
separan. En el formato \textbf{narrativo}, la historia se escribe directamente desde el plan y las
promesas: un redactor escribe cada capítulo, un crítico lo anota y un escritor lo corrige párrafo a
párrafo. Es la generación \textit{Top-Down} con LLMs de la Sección~\ref{sec:top-down}. En el formato
de \textbf{guion}, el resultado es el propio guion teatral. En el formato \textbf{simulado}, el
guion se representa y se narra; es el generador híbrido que propone esta tesis. Como cada etapa
guarda su resultado antes de dar paso a la siguiente, una historia se puede repetir desde cualquier
punto: escribir en prosa directa el plan de una función ya representada, o narrar de nuevo la misma
función desde otro punto de vista. Esta reutilización permite comparar los rasgos de personajes y trama y la preferencia de los lectores
entre una historia narrada a partir de una función y otra escrita directamente desde el mismo plan.

Un principio recorre todas las etapas: \textbf{el modelo propone y el código decide}. Cada vez que
un agente produce algo con estructura, sea un plan, un acto, un dosier de personaje o un turno, el
código lo valida contra invariantes explícitos antes de aceptarlo. Si no los cumple, los mensajes de
error se reinyectan en el \textit{prompt} y el agente lo intenta de nuevo; lo que admite una sola
corrección se normaliza en lugar de rechazarse. Los hechos de los que depende el resto del sistema,
como el orden de los eventos, quién presenció cada turno o qué guarda cada memoria, los calcula el
código y nunca se le preguntan al modelo. Es la lección de la Sección~\ref{sec:consistencia}: los
fallos de información de los LLMs se previenen mejor por construcción que pidiendo al modelo que no
los cometa. Si una etapa no esencial falla, como el reparto o la narración de un capítulo, el código
deriva un resultado de respaldo, la ejecución lo registra y la historia sigue.

\section{El plan validado}\label{sec:plan}

La petición es un texto libre o una ficha de la obra con su trama y su reparto. Un analista la
interpreta y elige uno de tres perfiles, Esencial, Desarrollada o Expansiva, que fijan la escala de
la historia. Los perfiles se describen al modelo con palabras y no con presupuestos: si la persona
pide una extensión, el código la traduce al perfil en el que encaja, y la cifra no llega a ningún \textit{prompt}.

La petición del usuario se transforma en una especificación estructurada que recoge las condiciones
de la historia y el perfil narrativo correspondiente. Cada agente recibe instrucciones para su tarea
y los datos pertinentes de las etapas anteriores: el constructor del mundo define lugares y objetos,
y el diseñador de personajes utiliza ese contexto para elaborar el reparto. Cuando está activa, la
guía narrativa recupera esqueletos de trama relacionados con la petición y los incorpora como
orientación estructural. El planificador recibe la especificación, el mundo, los personajes y la guía
disponible para proponer un plan que después valida el código.

El plan es la decisión central del autor. Se compone de capítulos, cada uno con su objetivo
dramático, su estado inicial, su punto de giro y su estado final, y de eventos. Cada evento declara
su propósito, su función dramática, su conflicto y su desenlace; las condiciones que lo hacen
posible y los cambios que produce, en lenguaje natural; los personajes, el lugar y los objetos que
intervienen, y, si es el caso, las preparaciones anteriores que paga. Las dependencias entre
eventos son causales o temporales.

Sea $E$ el conjunto de eventos y $D \subseteq E \times E$ el de dependencias, donde $(e_i, e_j) \in
D$ significa que $e_i$ debe ocurrir antes que $e_j$. El código acepta el plan solo si el grafo
$(E, D)$ es acíclico, los identificadores son únicos y los órdenes consecutivos, toda referencia a
un capítulo, personaje, lugar u objeto existe y cada capítulo tiene al menos dos eventos. Con el
grafo aceptado, el código calcula un orden topológico estable $\sigma$, que asigna a cada evento su
posición y es el único orden en que confía el resto del sistema; todo evento $e$ que paga una
preparación $e'$ debe cumplir $\sigma(e') < \sigma(e)$. Un plan que no valida se rechaza, sus
errores se reinyectan en un \textit{prompt} de reparación y se pide otro. Como DOC~\cite{yang2023} y
Agents' Room~\cite{huot2025}, el generador planifica antes de escribir; a diferencia de ellos, el
plan no es solo una guía para el modelo, sino un objeto que el código comprueba. PLOTTER~\cite{gu2026}
también acepta solo las ediciones que dejan su grafo causal acíclico y sus eventos alcanzables, pero
entre esas restricciones no está el orden entre una preparación y su pago. La validación como
grafo acíclico retoma el gestor de dependencias de Fuentes Rodríguez~\cite{fuentes2025}, que ya
representaba la trama como un grafo dirigido acíclico de eventos.

\subsection{Promesas, progreso y pago}\label{sec:promesas}

Piad Morffis~\cite{piad2026}, a partir de Sanderson, describe la trama como una tríada: la historia
hace \textbf{promesas} al principio, muestra un \textbf{progreso} hacia ellas y las \textbf{paga}
al final, con tantos pagos como promesas. La tríada es, en sus palabras, «esencialmente, un contrato
verificable», y el generador la trata así. Con el plan ya congelado, un agente traza el contrato de
promesas: para cada una, el evento que la abre, los que la hacen progresar y el que la paga, todos
citados por su identificador. El capítulo de cada momento no lo escribe el modelo, sino que se
deriva del evento, de modo que el contrato no puede inventar un momento que el plan no tenga.

El código valida el contrato sobre el orden $\sigma$. Si $a$ es la apertura de una promesa, $g$
cualquiera de sus progresos y $r$ su pago, debe cumplirse $\sigma(a) < \sigma(g) < \sigma(r)$; la
promesa principal se paga en el último capítulo, y el número de promesas cae dentro de la banda del
perfil. Validado, el contrato se convierte en obligaciones por capítulo para los agentes que
escriben: cambia cómo se escribe el plan, nunca lo que contiene. Es la propuesta de Yun et
al.~\cite{yun2026}, representar cada preparación y su pago de forma comprobable, llevada a los
eventos de un plan. El contrato se puede apagar, y ese es su brazo de control: una auditoría
registra qué promesas pagó la historia final.

\section{El guion teatral}\label{sec:guion}

Un plan dice qué eventos ocurren, pero no quién está en escena, dónde ni qué quiere cada uno en ese
momento. Eso lo decide el guion. Un dramaturgo escribe cada capítulo del plan como un acto dividido
en escenas. Cada escena se ancla a uno o varios eventos del plan, en su orden, y declara su lugar,
una acotación inicial con lo que vería el público, el reparto presente con el objetivo concreto que
cada personaje persigue en ella, y sus réplicas y acotaciones. El objetivo no se imprime en el guion,
pero es lo que después recibe el actor.

El código valida cada acto contra el plan congelado: rechaza los anclajes a eventos inexistentes o
desordenados y los lugares y repartos que no corresponden, y normaliza lo demás. Un crítico revisa el
acto y un escritor lo corrige. El guion se puede obtener de dos maneras, con el mismo validador: el
método \textbf{nativo}, en que el dramaturgo escribe los actos desde el plan, y el \textbf{adaptado},
en que se adapta a actos la prosa del formato narrativo. Ambos se implementaron para compararlos a
ciegas y quedarse con uno; esa comparación quedó fuera de los experimentos, y la configuración
de referencia usa el nativo, que es el que alimenta la función.

En el formato simulado, el guion es la partitura de la función, no su texto. Cada escena del guion
se convierte en una escena de la función, con un \textit{beat} por cada evento del plan que
representa; el término viene del teatro y, en la generación de historias, de Façade
(Sección~\ref{sec:hibridos}). Las réplicas que escribió el dramaturgo no las recibe nadie en la
función: solo sirven para medir después si los actores, aun sin verlas, acabaron recitándolas.

\section{La función simulada}\label{sec:funcion}

La función es la mitad \textit{Bottom-Up} del generador. Su propósito es que lo que dicen y hacen
los personajes salga de lo que cada uno sabe y quiere, y no de la visión del autor. Tiene tres
piezas: el reparto, que fija qué sabe cada personaje al empezar; los actores, cada uno con su
memoria, y el director, que vela por que cada escena alcance los eventos del plan.

\subsection{Reparto y compuertas de conocimiento}\label{sec:reparto}

Antes de la primera escena, un director de reparto escribe un dosier por personaje: su cara pública,
sus tácticas habituales, las reglas que nunca rompe, sus relaciones, su emoción y su objetivo
iniciales y, sobre todo, lo que ya sabe al levantarse el telón, con la procedencia de cada hecho. Si
el personaje guarda un secreto, el dosier incluye la versión que cuenta en su lugar. Lo que no está
en ese conocimiento inicial, el personaje tiene que aprenderlo en escena.

La asimetría de información de la historia se declara en \textbf{compuertas de conocimiento}. Cada
compuerta es un hecho que unos personajes saben y otros no, con quién lo sabe antes de la primera
escena, quién lo sacará a la luz, de qué modo (confesión, deducción, descubrimiento, porque se lo
cuentan o porque lo oye) y en qué evento del plan pasa a ser sabido. El código hace cumplir una
distinción que el modelo confundía: quien descubre un hecho en escena no puede saberlo desde el
principio. La primera ejecución real le dio al detective la solución de partida. Cada actor recibe
solo los hechos que ya conoce, y únicamente el director ve todas las compuertas. Es la idea de
Sabre~\cite{ware2021}, en que cada personaje actúa a partir de lo que cree, impuesta sin un
planificador formal: un secreto se guarda porque el actor que no debe saberlo nunca lo recibe, no
porque se le pida al modelo que lo guarde.

\subsection{Los actores y su memoria}\label{sec:memoria}

Cada personaje lo interpreta un actor: una llamada al LLM por turno, con un contexto propio. El
actor nunca ve el guion ni el plan; recibe las circunstancias de la escena, su objetivo en ella, la
nota del director y lo que recupera de su memoria. Ni los identificadores ni los títulos de los
eventos ni las escenas futuras llegan a su contexto, y una prueba recorre una ejecución completa
para comprobarlo. En cada turno, el actor declara qué dice, qué gesto visible hace, qué piensa sin
decirlo, a quién se dirige, si habla en público o en susurro y qué táctica intenta, de una lista
cerrada (enfrentar, negar, mentir, confesar, ceder…). El código valida el turno, lo rechaza si se
sale de la ficción o del idioma de la historia y deja que el actor lo repare una vez.

La percepción la decide el código, no el actor. Un turno público lo perciben los personajes
presentes en la escena; un susurro, solo quien habla y sus destinatarios, y un pensamiento, solo
quien lo piensa. Cada personaje tiene un flujo de memoria propio en el que se escribe únicamente lo
que percibió. No hay un almacén común del que filtrar al leer: lo que un personaje no presenció
nunca se escribió en su memoria. Es la respuesta del generador a fallos como el del personaje que
revela lo que aún no podía saber~\cite{ahn2024} y a la dificultad de los LLMs para razonar quién
sabe qué~\cite{kim2023}: en lugar de pedirle al modelo que lo razone, el contexto de cada actor ya
contiene solo lo que su personaje sabe.

Para cada turno, el actor no recibe toda su memoria, sino los registros más pertinentes. La
recuperación sigue a los agentes generativos de Park et al.~\cite{park2023}, que combinan recencia,
importancia y relevancia, con dos diferencias: no usa \textit{embeddings} (representaciones
vectoriales del texto) ni llamadas al modelo, de modo que la misma ejecución puntúa igual en
cualquier máquina, y añade un término de compañía. La puntuación de un registro $m$, formado en la
escena $k_m$, al preparar un turno de la escena $k$ es
\begin{equation}
  s(m) = 1{,}0\,\mathrm{rel}(m) + 0{,}6\,\frac{1}{1 + 0{,}25\,(k - k_m)} + 0{,}5\,\mathrm{imp}(m)
  + 0{,}3\,\mathrm{comp}(m),
\end{equation}
donde $\mathrm{rel}(m) \in [0, 1]$ es el solapamiento léxico entre el registro y la situación de la
escena, en el que una palabra pesa más cuanto más rara es en el pasado del personaje;
$\mathrm{imp}(m) \in [0, 1]$, una importancia fijada por reglas para los turnos y propuesta por el
modelo para las reflexiones, y $\mathrm{comp}(m)$, la fracción de los personajes presentes que
intervinieron en él. Al actor le llegan los seis registros mejor puntuados y sus dos reflexiones más
recientes.

Al terminar cada escena, cada actor reflexiona sobre lo vivido, y su estado, con sus objetivos y su
postura hacia los demás, se sustituye en lugar de acumularse: un actor que guardara a la vez «es mi
aliada» y «me traicionó» interpretaría mal las dos. El historial queda en los registros. La memoria
propia tiene, además, su brazo de control. Con memoria compartida, cada turno público llega a la
memoria de todo el reparto, esté o no en escena, mientras que susurros y pensamientos siguen siendo
privados. Comparar ambas configuraciones sobre el mismo plan permite comprobar si la memoria propia reduce
las menciones a hechos que el personaje no podía conocer y si los lectores prefieren las escenas resultantes.

\subsection{El director}\label{sec:director}

El director es el componente con visión de autor dentro de la función. Para cada \textit{beat} ve
el evento del plan que hay que alcanzar, con su propósito, su conflicto y su desenlace, las
obligaciones del contrato de promesas y las compuertas que deben abrirse en él. Como el director de
The Virtual Storyteller~\cite{theune2003}, sugiere motivaciones y nunca dicta réplicas: abre cada
\textit{beat} indicando quién se mueve, y su nota le dice al actor qué buscar, no qué decir.

Cada tres turnos, el director lee el \textit{beat} cláusula por cláusula, y lo da por alcanzado
cuando cada cláusula tiene un turno que la muestra. Si en la segunda lectura sigue faltando algo,
señala quién tiene que cambiar y por qué. Cuando el \textit{beat} agota su presupuesto de turnos,
ocho por defecto, el mundo entrega la cláusula que falta con un \textbf{suceso del mundo}: algo
visible que ocurre en escena y entra al registro como un turno más, de modo que el registro sigue
siendo la única fuente de la historia. Dos actores reaccionan y una última lectura decide si el
\textit{beat} se alcanzó con ayuda o se forzó. En el último \textit{beat} de la obra, el suceso del
mundo puede ejecutar la consecuencia visible de una decisión que los personajes ya tomaron, pero
nunca decidir por ellos.

Este mecanismo garantiza que la función avance, no que los personajes lleguen solos a cada evento.
La función se representa por defecto en modo \textbf{fijo}, en el que los eventos del plan son hitos que hay que
alcanzar. El modo \textbf{adaptativo}, opcional, permite que un hito no representado revise las
escenas futuras, con una propuesta y una reparación por frontera entre escenas; si fallan, la función
termina con un cierre abierto, y el plan original nunca se modifica.

La función registra, además, medidas de diagnóstico, como cuánto reproducen los actores las réplicas
del guion que nunca vieron. Comprueban el mecanismo, no la calidad de la historia, que es objeto del
Capítulo~\ref{chapter:evaluation}.

\section{El narrador}\label{sec:narrador}

El narrador convierte el registro de la función en la historia, con un capítulo por cada capítulo
del plan. Como en StoryBox~\cite{chen2026a}, lo que hacen los agentes y quién lo cuenta están
separados, pero aquí el narrador tiene un límite estricto: \textbf{cura sin inventar}. Puede cortar,
fundir y reordenar lo ocurrido dentro de un capítulo, pero no añadir un suceso que no esté en el
registro, y donde el plan y el registro no coinciden, manda el registro. La historia cuenta lo que la
función produjo, no lo que el plan preveía.

Darle el registro tal cual no funciona, porque el narrador copia lo que recibe: con el registro
entero transcribía todas las réplicas, y con las réplicas recortadas copiaba los fragmentos. Por
eso, en la narración \textbf{resumida}, que es la predeterminada, recibe enteras solo las réplicas
clave de cada escena y los sucesos del mundo. El resto le llega en estilo indirecto y en pasado
(«el capitán insistió en el lastre»), junto con los gestos y la ambientación de cada escena, todo
ello resumido por otra llamada al modelo una vez por capítulo. El código descarta el resumen que
copia su réplica, el que comparte con ella cuatro palabras seguidas. Después, un revisor comprueba
cada capítulo contra el registro: solo lo bloquean hallazgos graves, como un suceso inventado o una
contradicción, y el hallazgo que no puede citar el turno en que se basa se descarta.

El punto de vista es modular: cada visión es un filtro determinista sobre el registro. La
\textbf{omnisciente}, la predeterminada, ve todos los turnos; la \textbf{limitada} y la de
\textbf{primera persona} siguen a un personaje y solo ven las escenas en que estuvo presente. Un
capítulo sin turnos visibles no se narra, porque el narrador tendría que inventarlo. La visión, el
personaje y el tono que pide la persona solo llegan a la narración, y una prueba comprueba que el
plan, el guion, el reparto y la función no cambian entre dos ejecuciones que solo difieren en ellos.
Así, la misma función se puede narrar con varias visiones o tonos, y las diferencias entre las
historias resultantes se deben solo al discurso.

\section{Mecanismos opcionales}\label{sec:opcionales}

WSE-bench~\cite{chen2026b} concluye que añadir estructura no mejora la coherencia de forma uniforme,
y este trabajo contrasta esa observación con la preferencia humana: que un mecanismo mejore medidas
de la prosa no basta para que los
lectores prefieran las historias que lo usan. Por eso el generador implementa más mecanismos de los
que usa su configuración de referencia, todos con una regla de diseño: apagados, dejan el generador
como estaba; en el inventario, una prueba automática comprueba que no cambia ningún
\textit{prompt}, esquema ni resultado intermedio. Cada uno es así un brazo experimental, y su
estado por defecto cambia solo por una lectura humana a ciegas, no por una medida automática. La Tabla~\ref{tab:mecanismos} resume los mecanismos con
alternativa, su etapa y su estado en la configuración de referencia.

\begin{table}[htb]
  \centering
  \caption{Estado de los mecanismos en la configuración de referencia}
  \label{tab:mecanismos}
  \small
  \begin{tabular}{@{}>{\raggedright\arraybackslash}p{5.4cm}>{\raggedright\arraybackslash}p{2.8cm}%
      >{\raggedright\arraybackslash}p{2.0cm}>{\raggedright\arraybackslash}p{3.6cm}@{}}
    \toprule
    Mecanismo & Etapa & Referencia & Alternativa \\
    \midrule
    Guía de esqueletos de trama & Plan & encendida & apagada o compositiva \\
    Contrato de promesas & Plan & encendido & apagado (control) \\
    Método del guion & Guion & nativo & adaptado de la prosa \\
    Memoria de los actores & Función & propia & compartida \\
    Modo de la función & Función & fijo & adaptativo \\
    Cláusulas por tipo & Función & apagadas & encendidas \\
    Inventario de objetos & Función & apagado & encendido \\
    Método de narración & Narración & resumida & por selección \\
    Visión & Narración & omnisciente & limitada o primera persona \\
    Fichas de escena y crítico con cita & Escritura narrativa & apagados & encendidos \\
    Editor por fragmentos y ejemplos de estilo & Escritura y narración & apagados & encendidos \\
    \bottomrule
  \end{tabular}
\end{table}

El más elaborado es el \textbf{inventario de objetos}, que aplica a las cosas el principio de la
Sección~\ref{sec:arquitectura}: el actor declara un movimiento, un verbo y el nombre de un objeto, y
el código lo valida contra la verdad que él mismo guarda. Cada uno de los seis verbos tiene una
regla (usar, dar, soltar, esconder o mostrar exigen tener el objeto; tomar, que esté a la vista), y
un objeto solo lo percibe quien lo ve. Su validación mostró también lo que el árbitro no cubre: el
suceso del mundo con que el director cierra un \textit{beat} es texto libre, y en una ocasión hizo
caer un sello de las manos de un personaje que, según el inventario, no lo tenía.

Los demás mecanismos atacan defectos de la prosa que el código puede medir. Las \textbf{fichas de
escena} convierten cada evento en una ficha (lugar, deseos, obstáculo, giro) antes de redactar; el
\textbf{crítico con cita} nombra el problema principal de cada capítulo con una cita literal; el
\textbf{editor por fragmentos} reescribe solo los párrafos que el código marca, y los
\textbf{ejemplos de estilo} añaden dos párrafos de muestra al contrato de estilo. La narración
\textbf{por selección} elige primero los momentos del registro y escribe desde ellos, y las
\textbf{cláusulas por tipo} separan lo que deciden los personajes de lo que hace el mundo.

La configuración de referencia conserva estos mecanismos como opciones que pueden activarse para
comparar su efecto. Su evaluación combina medidas de la prosa con lecturas humanas, porque reducir
un defecto medido no garantiza una preferencia por el texto resultante. El
Capítulo~\ref{chapter:experiments} presenta las comparaciones documentadas y el alcance de las
lecturas adicionales comunicadas por el autor.

Este capítulo deja descrito un generador que separa qué pasa, cómo pasa y cómo se cuenta, y que
apoya cada separación en el código: un plan que se valida como grafo, un contrato de promesas
anclado a sus eventos, actores que solo recuerdan lo que percibieron y un narrador que no puede
inventar. Deja también fijada su configuración, con los mecanismos que no se ganaron su
complejidad apagados pero disponibles. El Capítulo~\ref{chapter:evaluation} describe cómo se evalúan
las historias que produce.

\chapter{Metodología de evaluación}\label{chapter:evaluation}
% R04 en ../write_thesis_roadmap.md

\chapter{Implementación}\label{chapter:implementation}
% R05 en ../write_thesis_roadmap.md. Fuentes: AGENTS.md, docs/architecture/runtime.md e
% interfaces.md, los README de packages/* y apps/*, y el código el 2026-11-01: versiones de los
% pyproject.toml y de version.py (Studio es 0.3.0 desde el 2026-10-27; el «Estado actual» de
% planning/roadmap.md decía 0.2.2 y se corrigió), comandos de asg_core/help.py, temperaturas y
% reintentos de runtime/provider.py, constantes de features.py, study.py y judge.py. Líneas de
% código: líneas de los .py versionados bajo src/ de cada paquete (git ls-files), con blancos y
% comentarios, el 2026-11-01. Pruebas: .\run-tests.ps1 el 2026-11-01, 907 pasan y 2 se omiten
% (las de test_gemini_live.py). La cifra de los 503 es de docs/architecture/runtime.md.
% Referencias nuevas, verificadas en la web el 2026-11-01: pydantic2026, google2026,
% virtanen2020, pedregosa2011, ptb2026, fastapi2026. \codigo está definido en Thesis.tex.

Este capítulo describe cómo se implementaron el generador del Capítulo~\ref{chapter:generator} y
las herramientas de evaluación del Capítulo~\ref{chapter:evaluation}. A diferencia de esos
capítulos, aquí se nombran los paquetes, los módulos y los comandos del código. La descripción
corresponde al estado del repositorio el 1 de noviembre de 2026. Primero se presenta la
\textbf{Organización del código} en cinco paquetes (Sección~\ref{sec:monorepo}). Después se
describen los \textbf{Proveedores de modelos y la cuota} (Sección~\ref{sec:proveedores}), los
\textbf{Artefactos y contratos versionados} (Sección~\ref{sec:artefactos}), las \textbf{Herramientas
de evaluación} (Sección~\ref{sec:herramientas}) y las \textbf{Interfaces}, un bot de Telegram y una
aplicación web local (Sección~\ref{sec:interfaces}). El capítulo cierra con las \textbf{Pruebas y la
reproducibilidad} (Sección~\ref{sec:pruebas}).

\section{Organización del código}\label{sec:monorepo}

El sistema está escrito en Python, a partir de la versión 3.12, y se organiza como un
monorepositorio: un solo repositorio con cinco paquetes que se instalan por separado y que un script
de preparación instala juntos en modo editable. La Tabla~\ref{tab:paquetes} los resume. El generador
es el paquete \texttt{asg-stagecraft}, y la interfaz gráfica, llamada StageCraft, el paquete
\texttt{asg-studio}. Las dependencias van en un solo sentido: \texttt{asg-core} no depende de
ningún otro, \texttt{asg-evaluation} depende solo de él, el generador depende de ambos y las dos
aplicaciones, de los tres.

\begin{table}[htb]
  \centering
  \caption{Paquetes del repositorio el 1 de noviembre de 2026}
  \label{tab:paquetes}
  \small
  \begin{tabular}{@{}>{\raggedright\arraybackslash}p{3.0cm}>{\raggedright\arraybackslash}p{1.3cm}%
      >{\raggedright\arraybackslash}p{7.2cm}r@{}}
    \toprule
    Paquete & Versión & Responsabilidad & Líneas \\
    \midrule
    \texttt{asg-core} & 0.8.1 & escritura atómica, medidas de la prosa, narración en audio,
      ayuda de los comandos & 1.065 \\
    \texttt{asg-evaluation} & 2.1.6 & rasgos con cita, estudio por pares, Bradley-Terry, juez y
      ranking & 3.467 \\
    \texttt{asg-stagecraft} & 7.22.1 & el generador: agentes, validadores, función, narrador,
      proveedores y cuota & 23.036 \\
    \texttt{asg-telegram} & 6.0.3 & bot para crear historias y para el estudio a ciegas & 1.540 \\
    \texttt{asg-studio} & 0.3.0 & StageCraft, interfaz web local y lectura ciega & 1.483 \\
    \bottomrule
  \end{tabular}
\end{table}

La restricción más importante es que \texttt{asg-evaluation} nunca importa el generador: lee del
disco los artefactos que este deja, en JSON, Markdown y SQLite. Así la evaluación mide del mismo
modo las historias de la versión congelada, las de versiones anteriores y las líneas base, que no
pasan por el generador, y ningún cambio en él puede alterar en silencio cómo se le evalúa. De ahí
dos consecuencias. Las medidas de la prosa que necesitan los dos lados, como la proporción de
párrafos con diálogo, viven en \texttt{asg-core}. Y el extractor de rasgos, que pertenece a la
evaluación pero necesita un LLM, recibe el modelo desde fuera: el comando \texttt{asg-features} es
del generador, construye el proveedor y se lo entrega a la función de extracción, que solo conoce su
interfaz.

Dentro del generador, los agentes del Capítulo~\ref{chapter:generator} están en \texttt{agents/},
cada uno con su \textit{prompt} y el esquema de su respuesta, y cada validación determinista vive
aparte, en un módulo puro que no llama al modelo: el grafo del plan en \codigo{planning/graph.py},
el contrato de promesas en \codigo{planning/promises.py}, cada acto del guion en
\codigo{script/validation.py}, el reparto y sus compuertas en \codigo{stage/casting.py} y cada turno
propuesto en \codigo{stage/validation.py}. También son puras la percepción
(\codigo{stage/perception.py}), la memoria y su recuperación (\codigo{stage/memory.py}), la decisión
de quién habla y cuándo se consulta al director (\codigo{stage/policy.py}), las visiones del narrador
(\codigo{stage/voices.py}) y el árbitro del inventario (\codigo{stage/inventory.py}). Esa separación
es la forma concreta de «el modelo propone y el código decide» (Sección~\ref{sec:arquitectura}): lo
que decide el código se puede probar sin un modelo. Las opciones de una ejecución se validan en un
solo lugar, \texttt{options.py}, en el que confluyen la línea de comandos, las variables de entorno y
las dos interfaces.

Todo el sistema se maneja desde la línea de comandos. \texttt{asg help} enumera los dieciséis
comandos, todos con el prefijo \texttt{asg-}, y la Tabla~\ref{tab:comandos} los agrupa por tarea.

\begin{table}[htb]
  \centering
  \caption{Comandos del sistema, agrupados por tarea}
  \label{tab:comandos}
  \small
  \begin{tabular}{@{}>{\raggedright\arraybackslash}p{2.2cm}>{\raggedright\arraybackslash}p{4.4cm}%
      >{\raggedright\arraybackslash}p{7.2cm}@{}}
    \toprule
    Tarea & Comandos & Qué hacen \\
    \midrule
    Generar & \texttt{asg-generate}, \texttt{asg-baseline}, \texttt{asg-matrix} & una historia,
      una línea base de un solo \textit{prompt} o una matriz de experimentos en serie \\
    Operar & \texttt{asg-budget}, \texttt{asg-recover}, \texttt{asg-recompute} & cuota restante,
      ejecuciones atascadas y medidas recalculadas de una función \\
    Medir & \texttt{asg-features}, \texttt{asg-audit-features}, \texttt{asg-report-features},
      \texttt{asg-corpus} & extracción y auditoría de rasgos, tabla de rasgos y análisis
      horizontal, inventario del corpus elegible \\
    Preferir y juzgar & \texttt{asg-study}, \texttt{asg-preferences}, \texttt{asg-judge},
      \texttt{asg-rank} & estudio por pares, Bradley-Terry, juez aprendido y ranking con las
      diez mejores \\
    Interfaces & \texttt{asg-telegram}, \texttt{asg-studio} & el bot y StageCraft \\
    \bottomrule
  \end{tabular}
\end{table}

\section{Proveedores de modelos y cuota}\label{sec:proveedores}

\texttt{LanguageModelProvider}, la interfaz por la que pasan todas las llamadas a un LLM, tiene dos
operaciones: generar texto libre y generar una respuesta estructurada según un esquema de
Pydantic~\cite{pydantic2026}. Ningún agente sabe qué modelo ni qué servicio le responde. La
implementación principal usa la API de Gemini~\cite{google2026}; en la configuración de ejemplo del
repositorio, el modelo principal es \texttt{gemini-3.5-flash-lite} y los turnos y reflexiones de los
actores van a \texttt{gemini-3.1-flash-lite}, de modo que una función reparte su gasto entre dos
cupos. El reparto lo decide un enrutador según la etapa y el agente de cada llamada. El director se
queda en el modelo principal, porque un juicio flojo sobre si un \textit{beat} se alcanzó termina
forzándolo.

Pydantic es la autoridad local sobre cada respuesta estructurada: un campo de más o de un tipo
equivocado la invalida. Los errores de validación se reinyectan en el \textit{prompt} y el modelo
tiene un segundo intento; si vuelve a fallar, la llamada termina en un error que los guarda. Es el
principio de la Sección~\ref{sec:arquitectura} en su nivel más bajo, debajo de las validaciones
propias de cada etapa. La temperatura tampoco la elige cada agente: sale de cinco perfiles con
nombre, que son extracción (0,15), revisión (0,2), reescritura (0,35), planificación (0,5) y prosa
(0,9).

La restricción práctica que más pesó en esta tesis fue la cuota. El nivel gratuito de la API limita
las peticiones diarias por proyecto y por modelo, del orden de 500, una cifra inferida de la
ejecución que la agotó y no confirmada por el servicio. El proveedor la administra en tres capas.
Los limitadores de peticiones por minuto se comparten en todo el proceso, con una clave por
credencial y modelo, de modo que dos generaciones simultáneas respetan un único presupuesto. Un
libro de cuota en SQLite, compartido entre procesos, apunta cada intento; \texttt{asg-budget} lo
muestra, y \texttt{asg-generate} se niega a empezar una historia si lo que queda del día no alcanza
para la media de su formato. Por último, un servicio saturado (error 503) se reintenta en el mismo
modelo tras 10 segundos y, al sexto intento, la llamada pasa al siguiente modelo, si lo hay. La
paciencia es corta porque cada 503 cuenta como una petición del cupo diario: el 13 de octubre de
2026, 301 de las 500 peticiones a \texttt{gemini-3.5-flash-lite} fueron 503.

Una cadena opcional de proveedores, vacía por defecto, pasa la llamada a Groq o a Mistral cuando
Gemini agota su cupo; la ejecución que cambia de modelo lo registra, y la evaluación no la empareja
con una de un solo modelo. Un fallo de configuración o de cuota aborta la ejecución; el de una etapa
no esencial, incluida la narración en audio, queda como aviso y la historia sigue.

\section{Artefactos y contratos versionados}\label{sec:artefactos}

Cada ejecución escribe una carpeta propia dentro de \codigo{Stories/Stagecraft/}, nombrada con la
fecha y el título, con los resultados de cada etapa: la petición, el mundo, los personajes, el plan,
el contrato de promesas y, según el formato, el guion, los dosieres del reparto, la memoria de cada
personaje, el registro de la función (\texttt{performance.json}), la narración y la historia
(\texttt{story.md}). La escritura es atómica: cada archivo se escribe primero en un temporal que
después reemplaza al definitivo, de modo que una ejecución interrumpida nunca deja un archivo a
medias. Un manifiesto guarda el SHA-256 y el tamaño de cada artefacto.

\texttt{metadata.json} registra el estado, las etapas completadas, los avisos y los modelos que
respondieron, y \texttt{llm\_calls.jsonl}, cada llamada con su agente, su modelo, sus
\textit{tokens} y sus reintentos. Los artefactos guardan además los bloques exactos de
\textit{prompt} que se inyectaron, como las obligaciones de cada capítulo derivadas del contrato de
promesas, de modo que una ejecución terminada se audita sin volver a derivar qué se le dijo a cada
agente. Y como cada etapa queda en disco, una ejecución nueva puede partir de otra:
\codigo{--plan-from} escribe en prosa directa el plan de otra ejecución, y \codigo{--narrate-from}
vuelve a narrar su función con otra visión o con otro tono. Con ellas se derivan las celdas de la
matriz del análisis horizontal (Sección~\ref{sec:horizontal}), y la primera es la que contrasta H1.

El generador distingue su versión (7.22.1) de la del contrato de sus artefactos (7.22), y cada
ejecución registra ambas. Una ejecución solo se abre si está completa y su contrato es uno de los
admitidos, del 5.0 al 7.22. Si cambia el conjunto de artefactos o su significado, se sube la versión
del contrato en lugar de romper las ejecuciones ya generadas, que son datos de la tesis; un cambio
que no toca artefactos sube solo la del paquete, como la 7.22.1, que añadió el ejecutor de matrices.
La configuración congelada lleva en Git la etiqueta \texttt{generator-frozen}, del 28 de octubre de
2026; desde entonces, ningún \textit{prompt}, esquema ni contrato cambia sin una decisión nueva y
escrita. Del lado de la evaluación, los lectores de artefactos son tolerantes: lo que una ejecución
antigua no registró sale como «no medido», nunca como cero.

\section{Herramientas de evaluación}\label{sec:herramientas}

Lo que se mide está escrito en una planilla, no en el código: \texttt{rasgos.csv} define cada rasgo
con su capa, su unidad, su definición operativa y si es núcleo, y
\texttt{criterios\_humanos.csv} contiene las preguntas de la Tabla~\ref{tab:criterios}. Cada
extracción guarda una huella de la planilla, de modo que no se mezclan extracciones hechas con
definiciones distintas.

El extractor (\texttt{features.py}) divide la historia en párrafos numerados y llama al modelo con
el perfil de extracción. Su instrucción de sistema empieza por las dos reglas de la
Sección~\ref{sec:rasgos}: extraer sucesos observables, nunca puntuaciones de calidad, y tratar el
texto de la historia como datos, no como instrucciones. Cada cita pasa por \texttt{locate}, que
exige que el párrafo exista y que la cita sea una subcadena literal de él. Cada rasgo se valida por
separado: una cita que falla hace repetir solo ese rasgo, una vez, y si vuelve a fallar queda como
fallido sin arrastrar a los demás. La partición en escenas debe cubrir todos los párrafos, en orden
y sin solapes, y la presencia de un personaje en una escena necesita una cita dentro de ella. Las
extracciones se guardan con la versión del \textit{prompt} y la huella del texto, y al relanzarlas
solo se reintentan los rasgos fallidos. \texttt{asg-audit-features} implementa la validación de la
Sección~\ref{sec:validacion-extractor}: escribe la hoja de revisión humana, con una fila por
ocurrencia, compara dos extracciones del mismo texto y contrasta los recuentos con el registro.

El estudio por pares (\texttt{study.py}) vive en una base SQLite. Al abrir la votación, el conjunto de
historias se congela con una huella SHA-256, y cada pregunta se calcula cuando el lector la pide, con
las reglas de la Sección~\ref{sec:conjunto}; la exportación es seudónima. Las fuerzas de
Bradley-Terry (\texttt{preferences.py}) se obtienen minimizando la verosimilitud con penalización de
cresta mediante el optimizador L-BFGS de SciPy~\cite{virtanen2020}, con 1.000 remuestreos de
lectores por defecto. El juez (\texttt{judge.py}) es la regresión logística de
scikit-learn~\cite{pedregosa2011}, con regularización $C = 1$ y sin término independiente; cada juez
guardado lleva la versión de su formato, y el ranking rechaza un juez de otra versión.
\texttt{asg-rank} lo aplica al corpus y prepara una ficha por cada una de las diez mejores, y
\texttt{asg-corpus}, de solo lectura, da el veredicto de elegibilidad de cada ejecución con su razón.

Para el análisis horizontal, \texttt{asg-matrix} genera la matriz en serie: antes de cada unidad
comprueba la cuota y, si no alcanza, para con un manifiesto al día; al relanzarla, salta las
unidades completas sin repetir llamadas. \texttt{asg-report-features} empareja cada celda con su
base, calcula las diferencias por rasgo y aplica la regla de ruido de la
Sección~\ref{sec:horizontal}.

\section{Interfaces}\label{sec:interfaces}

El bot de Telegram (\texttt{asg-telegram}), construido con python-telegram-bot~\cite{ptb2026}, corre
en la computadora del investigador. Para crear, ofrece cuatro opciones: petición libre o guiada paso
a paso, con el formato narrativo o con el simulado. Las peticiones entran en una cola SQLite con una
sola petición activa por persona, se genera una historia cada vez y cada una llega como archivo,
como texto y como audio, narrado con edge-tts. Un único módulo del bot importa el generador. Para
comparar hay dos flujos. Con \texttt{/evaluar}, cada persona compara a ciegas sus propias dos
historias, una de cada formato, con las tres preguntas; son votos de prueba que no entran en el
análisis. Con \texttt{/estudio} se vota el estudio de la Sección~\ref{sec:conjunto}: historias de
otras personas, solo como texto, sin título, encabezados ni audio, y la pregunta con A, B y «No
puedo decidir». Un candado por lector evita que una doble pulsación registre dos votos.

StageCraft (\texttt{asg-studio}) es una aplicación web local, servida con FastAPI~\cite{fastapi2026}
y escrita en HTML, CSS y módulos de JavaScript, sin paso de compilación. Sus vistas permiten crear
una obra viendo el \textit{prompt} exacto que se enviará, leer las ejecuciones guardadas con su
registro, también con los ojos de una visión, y comparar de dos a cuatro ejecuciones eje por eje,
indicando si la comparación está emparejada. La vista de \textbf{lectura ciega} muestra pares de
historias anónimas, X e Y, y guarda un voto por par; las respuestas de la aplicación no llevan
identificador, título ni configuración, y la clave de cada par se guarda aparte. Con ella se hicieron
las dos lecturas del 27 de octubre de 2026 que decidieron los brazos de escritura y de narración
(Sección~\ref{sec:opcionales}). StageCraft nunca escribe en las carpetas de las historias, solo
acepta escrituras con una cabecera propia y desde la propia máquina, y pinta como texto, nunca como
HTML, lo que produce el modelo.

\section{Pruebas y reproducibilidad}\label{sec:pruebas}

El 1 de noviembre de 2026, la suite completa pasó 907 pruebas y omitió las 2 que llaman a la API
real, que solo se ejecutan a petición porque consumen cuota; tardó algo más de dos minutos. Ninguna
otra prueba llama a un modelo: un proveedor falso recibe una secuencia de respuestas estructuradas,
puede forzar fallos en llamadas concretas y registra cada contexto que se le envía. Sobre una
ejecución completa con ese proveedor, varias pruebas comprueban los contratos del
Capítulo~\ref{chapter:generator}. \codigo{test_no_actor_ever_sees_the_plan_or_a_future_scene} reúne
los contextos de todos los actores y comprueba que en ninguno aparece el identificador ni el título
de un evento, ni el de un capítulo. \codigo{test_a_character_only_remembers_what_it_witnessed}
comprueba que cada registro de memoria procede de un turno que el personaje presenció.
\codigo{test_without_the_inventory_no_object_reaches_a_prompt_or_an_artifact} comprueba que, con el
inventario apagado, no se escribe ningún artefacto suyo ni llega a un \textit{prompt} ningún bloque
o esquema de objetos: un brazo apagado deja el generador como estaba. Y
\codigo{test_the_voice_only_changes_the_narration} genera dos ejecuciones que solo difieren en la
visión, el personaje y el tono, y exige que el plan, el guion, el reparto y la función sean idénticos
byte a byte y que las llamadas anteriores a la narración coincidan.

La reproducibilidad tiene un límite que conviene declarar: el modelo no es determinista, y dos
ejecuciones con el mismo \textit{prompt} y la misma versión producen historias distintas. Por eso la
evaluación trata ese ruido como un dato (Sección~\ref{sec:horizontal}) en lugar de suponer que una
ejecución se repite. Lo que sí se reproduce es todo lo que decide el código: dadas las mismas
respuestas del modelo, la validación, la percepción, la recuperación de la memoria, que no usa
\textit{embeddings} ni llamadas al modelo, y las medidas dan el mismo resultado. También se
reproduce el análisis: como cada ejecución conserva sus artefactos con su huella, la tabla de
rasgos, las preferencias, el juez y el ranking se recalculan desde el disco sin volver a generar, el
sorteo del estudio se repite con su semilla y cada historia declara la versión del generador, del
contrato y de los modelos que la produjeron.

Este capítulo deja descrito un sistema en el que las decisiones del Capítulo~\ref{chapter:generator}
tienen un lugar comprobable: los validadores y la memoria son funciones puras con pruebas, la
evaluación no puede importar el generador y cada ejecución queda en disco con su versión y su
huella. La cuota del servicio, administrada pero no evitada, acotó lo que se pudo generar. El
Capítulo~\ref{chapter:experiments} presenta lo que se midió con estas herramientas.

\chapter{Experimentos y resultados}\label{chapter:experiments}
% Revisión 2026-11-07: declaraciones del autor en fix.md; artefactos y hashes en
% ../tablas/revision.json. Los informes del piloto anterior conservan su alcance histórico.
% Validaciones del generador: docs/resumen.md y Evaluations/estudio/lecturas/.

Este capítulo presenta las \textbf{Condiciones y el corpus} (Sección~\ref{sec:corpus}), las
\textbf{Validaciones del generador} (Sección~\ref{sec:validaciones}), la \textbf{Validación del
extractor} (Sección~\ref{sec:resultados-extractor}) y el \textbf{Análisis horizontal}
(Sección~\ref{sec:resultados-horizontal}). Después describe el \textbf{Estudio humano}
(Sección~\ref{sec:resultados-estudio}), el \textbf{Juez aprendido}
(Sección~\ref{sec:resultados-juez}) y la \textbf{Ordenación del corpus}
(Sección~\ref{sec:resultados-ranking}). Cierra con la \textbf{Respuesta a las hipótesis}
(Sección~\ref{sec:respuesta-hipotesis}). Se distingue entre los resultados reconstruibles desde los
artefactos y la información del estudio actualizada por el autor el 7 de noviembre de 2026; los
recuentos de participación proceden de esta última.

\section{Condiciones y corpus}\label{sec:corpus}

Todas las historias de este capítulo se generaron con la API de Gemini, sujeta a la cuota diaria
descrita en la Sección~\ref{sec:proveedores}, y con los dos modelos de la configuración de ejemplo:
\texttt{gemini-3.5-flash-lite} como modelo principal y \texttt{gemini-3.1-flash-lite} para los turnos
y las reflexiones de los actores, según los modelos que registra cada ejecución. El extractor de
rasgos usó \texttt{gemini-3.5-flash-lite}. Los contrastes de generación documentados mantienen la versión del
generador y la entrada: los cambios de narración, sobre la misma función, y los de
escritura, sobre el mismo plan (Sección~\ref{sec:artefactos}).

El corpus es el conjunto de ejecuciones guardadas en la copia del autor del repositorio, en la
carpeta \texttt{Stories/}, que no está bajo control de versiones (Sección~\ref{sec:pruebas}). El
inventario del 28
de octubre de 2026, hecho con \texttt{asg-corpus} antes de generar la matriz, contó 331 ejecuciones,
de las que 116 eran elegibles para el ranking: las completas y en prosa de Stagecraft 7 y las líneas
base. Se excluyeron 135 por ser de un generador anterior a Stagecraft 7, 72 por no haber terminado y
8 por tener un formato fuera del estudio en prosa, como los guiones; ninguna se borró ni se editó a
mano. El 1 de noviembre de 2026, con la matriz ya generada, el mismo inventario cuenta 346
ejecuciones y 130 elegibles: 70 simuladas, 48 narrativas y 12 líneas base. De ellas, 26 son de la
versión de referencia (13 simuladas, 4 narrativas y 9 líneas base); las demás, de versiones anteriores de
Stagecraft 7, son elegibles según ese inventario. Este inventario histórico no identifica por sí
solo el corpus de 122 historias destinado a la ordenación en el estudio actualizado
(Sección~\ref{sec:resultados-ranking}).

Los experimentos usan tres obras del catálogo de siete peticiones preparado para ellos: una fantasía
canónica con un caballero, una princesa y un dragón (obra 01), un misterio, «La última luz del faro»
(obra 03), y un drama, «El cine de los domingos» (obra 04); sus peticiones están en el
Anexo~\ref{app:obras}. Sus líneas base de un solo
\textit{prompt} se generaron el 28 de octubre de 2026 con la versión de referencia: tres por obra, con la
variante Esencial literal del catálogo, la misma petición que recibe el generador, sin plan ni
revisión. Las nueve terminaron con una llamada cada una, salvo una que repitió tras un fallo del
proveedor. Su extensión va de 834 a 978 palabras en la obra 01, de 1.149 a 1.725 en la 03 y de 1.097
a 1.722 en la 04. Las tres repeticiones de cada obra estaban pensadas para medir el ruido de la
propia línea base (Sección~\ref{sec:horizontal}).

\section{Validaciones del generador}\label{sec:validaciones}

Los mecanismos con alternativa se probaron en ejecuciones reales y en sesiones de lectura.
Las comprobaciones técnicas y las sesiones conservadas del 27 de octubre documentan la elección de
la configuración de referencia de la Tabla~\ref{tab:mecanismos}. El autor informa además de otras
lecturas realizadas por participantes del estudio; sus respuestas no se agregan a las sesiones
conservadas sin un registro que permita identificar cada comparación. A continuación se presentan
las mediciones fechadas que pueden reconstruirse.

\subsection{El modo de la función}\label{sec:val-modo}

El 30 de septiembre de 2026, con la versión 7.5, dos funciones se representaron desde el mismo plan
de «La sombra del volcán», con el mismo presupuesto de turnos. En modo fijo, la función completó 6
escenas y 48 turnos, y de sus 6 \textit{beats} alcanzó 5 y forzó 1. En modo adaptativo terminó con un
cierre abierto tras 2 escenas y 16 turnos, con 1 de sus 2 \textit{beats} alcanzado. Ninguna cumplió
las dos promesas del plan. El modo fijo quedó como predeterminado y el adaptativo, como opción sin
lectura a ciegas. Son dos observaciones de una sola premisa: muestran que el modo adaptativo puede
cerrar la obra antes de tiempo, no cuál de los dos da mejores historias.

En esa validación se auditaron literalmente 62 contextos enviados a los actores, y ninguno contenía un
pensamiento de otro personaje ni un susurro que el actor no hubiera presenciado. Es la comprobación,
en ejecuciones reales, de lo que una prueba de contrato comprueba con el proveedor falso
(Sección~\ref{sec:pruebas}). No se conserva una tabla de esa auditoría, solo su resultado en
\texttt{docs/resumen.md}. Revisó dos categorías que el diseño ya mantiene privadas por construcción,
y con la versión 7.5, anterior a la referencia experimental. Para el contraste entre memorias es evidencia indirecta: es coherente con que la
memoria propia aísle a los personajes de esas dos fugas, pero no muestra que mejore las escenas.

\subsection{Narración y escritura}\label{sec:val-brazos}

El 27 de octubre de 2026, con la versión 7.21, se conservaron dos sesiones de lectura a ciegas de los brazos de
narración y de escritura de la Sección~\ref{sec:opcionales}. En la primera, cinco funciones ya
representadas se narraron de nuevo dos veces cada una: con la narración resumida (brazo A) y con la
narración por selección, el editor por fragmentos y los ejemplos de estilo (brazo B). En la segunda,
dos planes nuevos se escribieron en prosa directa dos veces por brazo: sin brazos (A) y con las
fichas de escena, el crítico con cita, el editor y los ejemplos (B). El autor leyó los pares en la
vista de lectura ciega de StageCraft (Sección~\ref{sec:interfaces}), sin saber qué brazo había
producido cada historia. La Tabla~\ref{tab:lecturas} resume lo que cambió en las medidas y lo que
eligió el lector.

\begin{table}[htb]
  \centering
  \caption{Sesiones documentadas de narración y escritura (autor, 2026-10-27)}
  \label{tab:lecturas}
  \small
  \begin{tabular}{@{}>{\raggedright\arraybackslash}p{2.2cm}>{\raggedright\arraybackslash}p{1.0cm}%
      >{\raggedright\arraybackslash}p{8.0cm}>{\raggedright\arraybackslash}p{2.2cm}@{}}
    \toprule
    Lectura & Pares & Medidas, brazo A frente a brazo B & Elección \\
    \midrule
    Narración de la función & 5 & réplicas con verbo de habla: 30--62~\% frente a 3--22~\%; párrafos
      con más de dos réplicas indirectas: de 1 a 6 por historia frente a ninguno en cuatro pares y 2
      en el quinto & A 2, B 2, ninguna 1 \\
    Escritura narrativa & 4 & proporción de réplicas con coletilla emocional: 0,00--0,21 frente a
      0,00; proporción de párrafos marcados por el código: 0,25--0,68 frente a 0,00--0,10; palabras:
      1.581--1.827 frente a 1.197--1.503 & A 3, B 1 \\
    \bottomrule
  \end{tabular}
\end{table}

En las dos lecturas, el brazo B corrigió los defectos medidos, y en ninguna ganó. En la narración, B
redujo los párrafos que encadenan réplicas en estilo indirecto, como un acta, y el verbo de habla en
casi cada línea, y empató con A. Estas dos cifras son aproximadas: salen de expresiones regulares y no
del código del repositorio. \texttt{docs/resumen.md} registra además dos medidas en contra de B. La
proporción de réplicas del registro que sobreviven en la narración, que calcula el código del
generador, fue de 0,73 a 0,85 frente a 0,48 a 0,67 de A. Las réplicas directas por ejecución fueron
de 18 a 68 frente a 12 a 35. B conservó, por tanto, más del registro tal cual, que es el riesgo de la
Sección~\ref{sec:narrador}. En la escritura, los ocho borradores ya
cumplían, con o sin brazos, los cuatro umbrales con que el código mide la prosa (nombres repetidos,
coletillas emocionales, gerundios en la coletilla y símiles); B redujo las coletillas a casi cero y
acortó el texto, y el lector prefirió A en tres de los cuatro pares. Con un empate se mantuvo la
narración más barata, y en la escritura, la preferida: todos los brazos quedaron apagados, como
opciones.

Al día siguiente, dos ejecuciones con la versión 7.21 comprobaron, sin corregir nada, dos defectos
conocidos. El plan de una petición con ficha detallada («Las islas del cielo», perfil Expansiva)
dejó fuera dos cosas que pedía la ficha, la división de la tripulación y quién paga el coste de la
decisión final, y el revisor del plan lo aprobó sin notas: no comprueba cada giro pedido, una
limitación observada en una sola ejecución. En cambio, un guion de la obra 03 salió con nombres,
papeles y reparto en español, sin una palabra en inglés.

\subsection{La configuración de referencia}\label{sec:val-congelada}

El 28 de octubre de 2026, la versión 7.22 se validó con tres ejecuciones de la obra 03, perfil
Esencial, con los valores predeterminados y sin corregir nada: una narrativa, un guion nativo y una
simulada. Las tres terminaron completas y con nombres en español. La narrativa hizo 22 llamadas y el guion, 18. La simulada hizo 236:
133 al modelo principal y 103 al de la función, 9 de ellas intentos fallidos que se repitieron, con
unos 570.000 \textit{tokens}. En la simulada se cumplieron todas las comprobaciones fijadas: los tres
capítulos salieron del narrador y no de un respaldo automático, ninguno de los 47 pensamientos de los
actores salió en inglés y ninguna línea de la historia quedó fuera de la ficción. Quedaron tres
defectos menores fuera de esos criterios en los textos conservados: una contracción omitida
(«de el faro») en la narrativa, un personaje de la ficha que el reparto del guion perdió y unos
títulos sin tilde en la simulada. Ese día se estableció la configuración de referencia para los experimentos, sin que ello
supusiera detener el desarrollo del generador.

La limitación conocida está en los \textit{beats}: la función forzó 3 de sus 6, en tres escenas
distintas. Es el precio del mecanismo de la Sección~\ref{sec:director}: el suceso del mundo garantiza
que la función avance, no que los personajes lleguen solos a cada evento. Las validaciones anteriores
no fijan la proporción: la de la versión 7.5 forzó 1 de 6 (Sección~\ref{sec:val-modo}), y
\texttt{docs/resumen.md} registra 5 de 10 en una de la 7.12, cuyos datos no se reproducen en esta
tesis. Con una ejecución, la cifra no separa el efecto del ruido.

Las sesiones documentadas muestran que reducir defectos medidos de la prosa no garantiza que el
lector prefiera el texto resultante: el brazo modificado empató en narración y perdió tres de los
cuatro pares de escritura. Esta observación es compatible con la hipótesis sobre los mecanismos
opcionales. Su alcance corresponde a esas sesiones; las lecturas adicionales requieren su propio
desglose para estimar cuánto se reproduce el resultado entre participantes.

\section{Validación del extractor}\label{sec:resultados-extractor}

La validación registrada el 29 de octubre de 2026 reúne doce historias: cinco simuladas con su
registro, cuatro narrativas y tres líneas base, de las versiones 7.8 a 7.22, con 21.471 palabras.
El modelo extractor fue \texttt{gemini-3.5-flash-lite}; en seis historias se repitió la extracción.
La Tabla~\ref{tab:extractor} reorganiza las hojas de revisión y el resumen de esas repeticiones.
Es una reconstrucción de los registros conservados, no una nueva revisión humana de las historias.

La precisión se calcula como ocurrencias correctas entre ocurrencias con veredicto correcto o
incorrecto. Las valoraciones sin resolución se mantienen fuera de ese denominador, conforme al
procedimiento de la Sección~\ref{sec:validacion-extractor}; sus recuentos se muestran en la
Tabla~\ref{tab:extractor-detalle} del Anexo~\ref{app:tablas}. Por ello, una precisión de 1,00 describe
las citas resueltas, no garantiza que se haya contado todo ni equivale a una validación completa.
Los campos de omisiones de las hojas están vacíos y no permiten calcular cobertura.

\begin{table}[p]
  \centering
  \caption{Validación registrada del extractor: citas y repetición (12 historias, 2026-10-29)}
  \label{tab:extractor}
  \small
  \setlength{\tabcolsep}{5pt}
  % Revisión documental 2026-11-07; fuentes en revision.json.
\begin{tabular}{@{}>{\raggedright\arraybackslash}p{5.1cm}rrrr@{}}
\toprule
Rasgo & \shortstack{Citas\\juzgadas} & Precisión & \shortstack{Pares\\comparables} & \shortstack{Diferencia\\media} \\
\midrule
Promesas de apertura & 5 & 1,00 & 5/6 & 0,00 \\
Señales de progreso & 15 & 1,00 & 4/6 & 0,00 \\
Preparación del pago final & 1 & 1,00 & 1/6 & 0,00 \\
Ciclos «sí, pero» & 5 & 1,00 & 2/6 & 0,00 \\
Ciclos «no, y además» & 5 & 1,00 & 2/6 & 0,25 \\
Proactividad del protagonista & 10 & 1,00 & 4/6 & 0,19 \\
Soluciones preparadas & 5 & 1,00 & 1/6 & 0,00 \\
Límites del mundo & 6 & 1,00 & 4/6 & 0,12 \\
Escenas con cambio & 7 & 1,00 & 2/6 & 0,00 \\
Decisiones con coste & 8 & 1,00 & 2/6 & 0,00 \\
Escenas con conflicto & 28 & 1,00 & 3/6 & 0,00 \\
Mentiras & 0 & --- & 1/6 & 0,00 \\
Secretos revelados & 3 & 1,00 & 3/6 & 0,00 \\
Complicaciones & 21 & 1,00 & 4/6 & 0,00 \\
Giros & 9 & 1,00 & 3/6 & 0,17 \\
Sucesos azarosos o externos & 6 & 1,00 & 4/6 & 0,18 \\
Usos de objetos & 21 & 1,00 & 4/6 & 0,22 \\
Contradicciones & 1 & 1,00 & 1/6 & 0,00 \\
\bottomrule
\end{tabular}

  \par\medskip
  \begin{minipage}{0.94\textwidth}
    \footnotesize
    Citas juzgadas: ocurrencias con veredicto resuelto; con menos de cinco, el apoyo es reducido.
    Pares comparables: historias con el rasgo medido en ambas extracciones, de seis repetidas.
    Diferencia media: promedio de $|a-b|/\max(|a|,|b|,1)$ sobre esos pares.
    El guion indica que no hay denominador; una diferencia nula no informa sobre las historias
    que no pudieron compararse.
  \end{minipage}
\end{table}
% Reconstrucción: verdict_summary(hojas/*.csv) y retest/resumen.json; hashes en revision.json.

La tabla presenta los dieciocho rasgos núcleo extraídos con citas. Los otros seis se calculan por
código: palabras, proporción de diálogo, repetición de trigramas, cuota del personaje más presente,
reparto de la participación y reparto del diálogo. Los tres últimos dependen de la segmentación o
las réplicas extraídas; por eso también se informa su repetibilidad en el anexo. Los cambios entre
los estados medido, no aplicable y fallido se separan allí de las diferencias numéricas.

La lista vigente contiene 24 rasgos núcleo. Los registros de validación corresponden al catálogo
anterior a esa selección; no se mezclan automáticamente con extracciones del catálogo vigente.
La presentación revisada conserva sus fechas y denominadores y permite identificar qué observaciones
necesitan una nueva revisión antes de atribuirles precisión o cobertura humana.

\section{Análisis horizontal}\label{sec:resultados-horizontal}

El conjunto presentado comprende \textbf{24 ejecuciones completadas}: cinco bases simuladas,
nueve líneas base de un solo \textit{prompt}, dos narrativas escritas desde el plan y ocho
renarraciones con cambios de tono o visión. Una unidad es una ejecución de una configuración para
una obra y una repetición determinadas; no es un rasgo ni una comparación. La
Tabla~\ref{tab:horizontal-conjunto} desglosa este conjunto y distingue la generación terminada de
la disponibilidad de una extracción núcleo completa y compatible.

\begin{table}[htb]
  \centering
  \caption{Conjunto horizontal y extracciones localizadas (revisión, 2026-11-07)}
  \label{tab:horizontal-conjunto}
  \small
  \setlength{\tabcolsep}{4pt}
  % Revisión documental 2026-11-07; fuentes en revision.json.
\begin{tabular}{@{}llrr@{}}
\toprule
Celda & Configuración & Ejecuciones & \shortstack{Extracción núcleo\\completa y compatible} \\
\midrule
M01 & Base simulada & 5 & 2 \\
M02 & Línea base de un prompt & 9 & 3 \\
M03 & Narrativa desde el plan & 2 & 0 \\
M04 & Tono sombrío & 2 & 2 \\
M05 & Tono humorístico & 2 & 0 \\
M06 & Visión limitada & 2 & 0 \\
M07 & Primera persona & 2 & 1 \\
 & Total & 24 & 8 \\
\bottomrule
\end{tabular}

\end{table}

Las líneas base y una base simulada proceden del 28 de octubre; las restantes ejecuciones se
completaron el 30 de octubre de 2026. El manifiesto, actualizado el 31 de octubre en UTC, contiene
45 posiciones históricas: las 24 completadas, una fallida y veinte pendientes. El
Anexo~\ref{app:tablas} conserva ese desglose. El total de 24 describe el conjunto final seleccionado
para este análisis, sin modificar retrospectivamente el diseño registrado.

La matriz reutiliza el plan o la función de una base: la escritura directa cambia el procedimiento
que produce la prosa, mientras que las variantes de tono y visión mantienen el registro de la
función. La Tabla~\ref{tab:coste} presenta el coste conservado por tipo de ejecución, incluidos los
intentos fallidos. Las nueve líneas base requirieron una llamada cada una, salvo el reintento ya
indicado en la Sección~\ref{sec:corpus}.

\begin{table}[htb]
  \centering
  \caption{Coste medido por unidad de la matriz (2026-10-28 y 2026-10-30)}
  \label{tab:coste}
  \small
  \begin{tabular}{@{}>{\raggedright\arraybackslash}p{4.6cm}>{\raggedleft\arraybackslash}p{2.0cm}%
      >{\raggedleft\arraybackslash}p{2.2cm}>{\raggedleft\arraybackslash}p{2.2cm}%
      >{\raggedleft\arraybackslash}p{2.2cm}@{}}
    \toprule
    Unidad & Ejecuciones & Modelo principal & Modelo de la función & \textit{Tokens} (miles) \\
    \midrule
    Base simulada & 5 & 49--133 & 49--115 & 299--570 \\
    Narrativa desde el plan de la base & 2 & 11--14 & 0 & 49--74 \\
    Renarración con otro tono o visión & 8 & 13--16 & 0 & 49--75 \\
    Base fallida (obra 04, repetición 2) & 1 & 0 & 198 & 320 \\
    \bottomrule
  \end{tabular}
\end{table}
% Llamadas: intentos por modelo en llm_calls.jsonl, incluidos los fallidos. Bases:
% 20261030-132707 (85 + 49), 20261030-234947 (49 + 115), 20261028-192554-el-misterio-del-faro
% (133 + 103), 20261030-140119 (102 + 74), 20261030-134427 (100 + 63). Fallida: 20261031-002920
% (198 al modelo de la función, 14 fallidas). Las carpetas llevan la hora UTC; en Cuba, todo es del
% 2026-10-30.

La comprobación documental del 7 de noviembre vincula las extracciones al texto mediante su huella,
el catálogo, el modelo y la versión del protocolo. Localizó ocho extracciones núcleo completas y
compatibles, cuatro compatibles con rasgos pendientes, seis de otro protocolo y seis unidades sin
archivo de extracción localizado en las carpetas revisadas. La Figura~\ref{fig:horizontal} muestra
cómo se reparten por celda: cinco de las seis extracciones de otro protocolo son de líneas base, y
cuatro de los seis archivos sin localizar corresponden a ejecuciones derivadas de la primera base de
la obra 01. Este inventario describe los archivos
disponibles; la declaración del autor de que las 24 fueron extraídas no sustituye los archivos
necesarios para reproducir sus medidas.

\begin{figure}[htb]
  \centering
  % Estado de la extracción núcleo por celda (fig:horizontal). Fuente:
% ../tablas/horizontal_revision.csv (columna estado_extraccion), comprobación del 2026-11-07.
% Recuentos a mano: completa 8, incompleta 4, otro protocolo 6, sin archivo 6; total 24.
\begin{tikzpicture}
  \begin{axis}[
      xbar stacked,
      width=0.72\textwidth,
      height=7.2cm,
      bar width=9pt,
      font=\footnotesize,
      xmin=0, xmax=10.2,
      xtick={0,1,...,9},
      xlabel={Ejecuciones},
      ytick={1,...,7},
      yticklabels={
        M01 Base simulada,
        M02 Línea base de un \textit{prompt},
        M03 Narrativa desde el plan,
        M04 Tono sombrío,
        M05 Tono humorístico,
        M06 Visión limitada,
        M07 Primera persona},
      y dir=reverse,
      ymin=0.4, ymax=7.6,
      axis x line*=bottom,
      axis y line*=left,
      xmajorgrids,
      major grid style={draw=black!12},
      tick style={draw=black!50},
      legend style={at={(0.5,1.03)}, anchor=south, legend columns=2, draw=none,
        font=\footnotesize, /tikz/every even column/.append style={column sep=8pt}},
      legend cell align=left,
    ]
    \addplot[fill=black!70, draw=black] coordinates
      {(2,1) (3,2) (0,3) (2,4) (0,5) (0,6) (1,7)};
    \addplot[fill=black!45, draw=black] coordinates
      {(0,1) (1,2) (1,3) (0,4) (1,5) (1,6) (0,7)};
    \addplot[fill=black!20, draw=black] coordinates
      {(1,1) (5,2) (0,3) (0,4) (0,5) (0,6) (0,7)};
    \addplot[fill=white, draw=black] coordinates
      {(2,1) (0,2) (1,3) (0,4) (1,5) (1,6) (1,7)};
    \legend{Completa y compatible (8), {Compatible, incompleta (4)},
      Otro protocolo (6), Sin archivo localizado (6)}
    % Total de ejecuciones de cada celda, al final de su barra.
    \node[anchor=west] at (axis cs:5,1) {5};
    \node[anchor=west] at (axis cs:9,2) {9};
    \node[anchor=west] at (axis cs:2,3) {2};
    \node[anchor=west] at (axis cs:2,4) {2};
    \node[anchor=west] at (axis cs:2,5) {2};
    \node[anchor=west] at (axis cs:2,6) {2};
    \node[anchor=west] at (axis cs:2,7) {2};
  \end{axis}
\end{tikzpicture}

  \caption{Extracción núcleo de las 24 ejecuciones, por celda (2026-11-07)}
  \label{fig:horizontal}
\end{figure}

Como ejemplo del procedimiento se conserva el análisis de las cinco variantes de la primera base
de la obra 03. La Tabla~\ref{tab:piloto} presenta los tres rasgos calculados directamente desde el
texto que resultaron comparables en el informe del 30 de octubre. Las diferencias corresponden a un
par por variante y no estiman el ruido entre repeticiones del conjunto completo.

\begin{table}[htb]
  \centering
  \caption{Diferencias con la base de la obra 03 (un par por variante, 2026-10-30)}
  \label{tab:piloto}
  \footnotesize
  \setlength{\tabcolsep}{4pt}
  \begin{tabular}{@{}lrrrl@{}}
\toprule
Celda & Palabras & Proporción de diálogo & Repetición de trigramas & Lectura \\
\midrule
Base (valores absolutos) & 1.560 & 0,40 & 0,12 & --- \\
Narrativa desde el plan & $-$176 & $+$0,06 & $-$0,04 & n insuficiente \\
Tono sombrío & $-$80 & $+$0,02 & $-$0,05 & n insuficiente \\
Tono humorístico & $-$75 & $-$0,08 & $-$0,04 & n insuficiente \\
Visión limitada & $-$311 & $-$0,05 & $-$0,04 & n insuficiente \\
Primera persona & $-$157 & $-$0,02 & $-$0,06 & n insuficiente \\
\bottomrule
\end{tabular}

\end{table}

El contraste entre simulación y escritura directa dispone de dos ejecuciones narrativas vinculadas
a sus planes fuente. Para interpretar la magnitud de sus diferencias según la regla de la
Sección~\ref{sec:horizontal}, deben acompañarse de las extracciones compatibles y del ruido entre
repeticiones. El autor comunica que se midió ese ruido con un resultado favorable, pero el informe
numérico correspondiente no está identificado en los registros revisados; no se le asigna un valor
ni se extiende a las 24 ejecuciones el resultado del ejemplo de la obra 03.

\section{El estudio humano}\label{sec:resultados-estudio}

Según la información del autor actualizada el 7 de noviembre de 2026, participaron
\textbf{ocho lectores}, incluido el propio autor, durante octubre y hasta el 1 de noviembre.
Se seleccionaron por conveniencia entre colegas graduados de Ciencia de la Computación. Tres
indicaron frecuencia de lectura alta, tres media y dos baja. Estas categorías describen hábitos
de lectura; no proceden de una escala validada ni representan diversidad de formación académica.
La Tabla~\ref{tab:lectores} del Anexo~\ref{app:lectores} recoge los perfiles anónimos.

Todos leyeron los mismos \textbf{diez textos distintos}, distribuidos en \textbf{cinco pares},
mediante el bot de Telegram. El método que produjo cada texto permaneció oculto durante la
comparación; conocer que las historias eran generadas no revelaba qué configuración correspondía
a A o B. El autor conocía el funcionamiento del sistema y se identifica como participante para
separar sus respuestas en el análisis. Los pares se presentaron en el mismo orden; el registro
actualizado no precisa si la posición A/B se mantuvo o se alternó entre lectores.

En cada par se respondieron las tres preguntas de la Tabla~\ref{tab:criterios}: creatividad,
desarrollo de personajes e interés de la trama. Los botones ofrecían «Historia A», «Historia B» y
«No puedo decidir»; la última opción se registra como abstención y no distingue entre empate y
falta de una preferencia. El autor confirma que los ocho completaron las preguntas. La
Tabla~\ref{tab:participacion} expresa la participación resultante: 40 comparaciones lector-par y
120 respuestas, incluidas las posibles abstenciones. Las tres respuestas sobre un par no constituyen
tres comparaciones independientes.

\begin{table}[htb]
  \centering
  \caption{Participación declarada del estudio (octubre--1 de noviembre de 2026)}
  \label{tab:participacion}
  \small
  % Revisión documental 2026-11-07; fuentes en revision.json.
\begin{tabular}{@{}lr@{}}
\toprule
Magnitud & Cantidad \\
\midrule
Participantes (incluido el autor) & 8 \\
Textos distintos & 10 \\
Pares comunes & 5 \\
Comparaciones lector-par & 40 \\
Respuestas por criterio & 40 \\
Respuestas totales (incluyen abstenciones) & 120 \\
\bottomrule
\end{tabular}

\end{table}

Las respuestas se transcribieron manualmente a partir de la actividad del bot. El registro
individual del conjunto de ocho lectores y la correspondencia entre sus diez textos y los
identificadores de ejecución no se han localizado en los archivos disponibles para esta revisión.
Por ello, se informa la participación declarada, pero no se reconstruyen elecciones individuales,
porcentajes de preferencia ni coeficientes de acuerdo. Las salidas conservadas de un piloto anterior,
con cinco lectores y seis textos, tienen otro alcance y no se utilizan como resultados de este
estudio. Sus identificadores de lector tampoco se equiparan con los del anexo actual.

Cinco pares disjuntos sobre diez textos forman cinco componentes en el grafo de comparaciones.
Incluso con todas las respuestas, esos pares permiten estudiar preferencias dentro de cada pareja,
pero no identificar un orden humano global mediante Bradley-Terry. Las lecturas adicionales,
realizadas por algunos de los mismos participantes, podrían aportar conexiones si se documentan sus
pares; no se presupone que lo hagan ni se incluyen en las 120 respuestas principales.

La muestra permite describir una aplicación de la metodología con distintos hábitos de lectura.
Para estimar preferencias y su incertidumbre se necesita recuperar las respuestas por par y criterio,
separar las del autor y remuestrear por lector, como prevé la Sección~\ref{sec:agregacion}.
Ampliar la muestra y diversificar su formación permitirá estudiar cuánto se reproducen los resultados.

\section{El juez aprendido}\label{sec:resultados-juez}

El juez relaciona los rasgos del texto con las preferencias humanas mediante la regresión por pares
de la Sección~\ref{sec:juez}. Deben distinguirse las historias que aportan votos al entrenamiento,
las reservadas para validación y el corpus al que después se aplica el modelo. El autor sitúa el
análisis actualizado el 1 de noviembre de 2026, con respuestas de ocho lectores, y comunica un rango
de acierto de \textbf{0,50 a 0,65}. La Tabla~\ref{tab:juez-actual} precisa el alcance documental de
ese resultado.

\begin{table}[htb]
  \centering
  \caption{Resultado comunicado del juez y alcance de su documentación}
  \label{tab:juez-actual}
  \small
  % Revisión documental 2026-11-07; fuentes en revision.json.
\begin{tabular}{@{}>{\raggedright\arraybackslash}p{5.2cm}>{\raggedright\arraybackslash}p{8.2cm}@{}}
\toprule
Dato & Alcance disponible \\
\midrule
Acierto comunicado & 0,50--0,65; sin desglose por criterio ni partición \\
Fecha comunicada & 1 de noviembre de 2026 \\
Procedencia & Declaración del autor, revisada el 7 de noviembre \\
Validación fuera de muestra & Modalidad y denominadores pendientes de documentar \\
Control de longitud & Sin resultado actualizado localizado \\
\bottomrule
\end{tabular}

\end{table}

No se dispone del desglose por criterio, los denominadores ni la partición que produjo ese rango.
Por tanto, no se presenta como una medida verificada fuera de muestra ni se compara con los
porcentajes de longitud sola del piloto anterior. Tampoco se calculan intervalos o significación a
partir de sus extremos. El informe del juez conservado del 3 de noviembre corresponde al piloto de
cinco lectores; sus predicciones y perturbaciones permanecen como resultados históricos separados.

La evaluación del juez debe conservar las predicciones de cada partición y el control de longitud
sobre los mismos votos. La interpretación requiere considerar la dependencia entre respuestas de
un lector y la falta de un ranking humano global en los cinco pares principales. Un rango de
acierto, por sí solo, no establece que el modelo generalice a nuevas historias o lectores.

\section{La ordenación del corpus}\label{sec:resultados-ranking}

El corpus destinado a la ordenación comprende \textbf{122 historias}, según la actualización del
autor. Este número describe los textos a los que se quiere aplicar el juez; no es el tamaño del
conjunto de entrenamiento ni demuestra que todos tengan extracción compatible o puntuación.
La lista de identificadores y la salida de ordenación correspondiente siguen pendientes de
incorporación al registro reproducible.

Los archivos históricos del 3 de noviembre describen otro alcance: nueve historias ordenadas de un
inventario de 130 elegibles, con seis textos del entrenamiento del piloto. Se conservan con su fecha
y sus fuentes, pero sus puestos y perfiles no se trasladan al corpus actualizado de 122 historias.
Sin la lista y las puntuaciones de este último no se seleccionan diez títulos ni se atribuyen a sus
primeras posiciones los rasgos de las historias del informe anterior.

Para interpretar la ordenación deben identificarse la versión del juez, las extracciones utilizadas,
las exclusiones y la pertenencia de cada texto al entrenamiento. La lectura de las primeras historias
ha de relacionar sus rasgos con pasajes concretos, sin convertir la puntuación del modelo en una
medida absoluta de calidad (Sección~\ref{sec:ranking}).

\section{Respuesta a las hipótesis}\label{sec:respuesta-hipotesis}

Cada hipótesis se responde según el contraste que propone y el alcance de las observaciones:

\begin{itemize}
  \item \textbf{Simulación frente a escritura directa desde el mismo plan (H1).} La matriz conserva
  dos narrativas derivadas de planes de bases simuladas y un informe de diferencias textuales para
  la obra 03. La asociación de los cinco pares humanos con esos contrastes requiere identificar
  los diez textos del estudio. Se conserva el contraste mecánico documentado; la preferencia
  específica por la simulación y su relación con el ruido deben establecerse con los registros
  correspondientes y una muestra ampliada de lectores.
  \item \textbf{Memoria propia frente a memoria compartida (H2).} El autor informa de una comparación
  realizada con resultado favorable a la memoria propia. La auditoría documentada de 62 contextos
  muestra aislamiento de pensamientos y susurros, pero no cuantifica el contraste entre ambas
  configuraciones. Los tres intentos con memoria compartida localizados en la versión 7.9 terminaron
  con fallo y no identifican el experimento comunicado. Para expresar su magnitud se necesitan los
  pares de ejecuciones, los recuentos de menciones desconocidas y las preferencias; no se deducen
  esos valores de la configuración del sistema.
  \item \textbf{Mejora mecánica y preferencia humana (H3).} Las sesiones conservadas de narración
  y escritura muestran mejoras en las medidas de prosa sin una victoria del brazo modificado:
  empate en narración y preferencia por el brazo sin esos mecanismos en tres de cuatro pares de
  escritura. Es evidencia descriptiva compatible con que una mejora mecánica no garantiza una
  preferencia humana. Las lecturas adicionales comunicadas deben documentarse para extender el
  alcance de esta observación; un empate o una diferencia no significativa no demuestran
  equivalencia entre configuraciones.
\end{itemize}

El capítulo documenta una configuración de referencia, 24 ejecuciones para el análisis horizontal y
un estudio de ocho lectores con 120 respuestas declaradas. Los instrumentos permiten relacionar
rasgos, comparaciones y configuraciones; su interpretación exige mantener la correspondencia entre
cada cifra y su fuente. El rango de acierto comunicado del juez y el corpus de 122 historias se
presentan con ese alcance, mientras que la generalización de las preferencias requiere ampliar la
muestra y completar su trazabilidad.


\backmatter

\begin{conclusions}
% R08 en ../write_thesis_roadmap.md. Una conclusión por objetivo específico, en el orden de
% Introduction.tex; ninguna afirma algo que no esté en los capítulos 1 a 5. Sin referencias nuevas.

Esta tesis se propuso desarrollar un generador híbrido de historias con LLMs, que combinara un plan
validado con la simulación de personajes con memoria propia, y una metodología para evaluarlo basada
en recuentos verificables y comparaciones humanas por pares a ciegas. A continuación se recoge lo
alcanzado en cada objetivo específico, en el orden en que se plantearon.

El análisis del estado del arte (Capítulo~\ref{chapter:state-of-the-art}) mostró que las tres
familias de generadores resuelven partes distintas del problema. Los dirigidos por la trama que se
revisaron garantizan la estructura y los compromisos, pero escriben a todos los personajes desde la
visión del autor; los que simulan personajes dan agencia, pero dejan la trama a la emergencia, y los
híbridos clásicos dependían de mundos formalizados a mano. En la literatura revisada, que no es
exhaustiva, no se encontró un generador con LLMs que reuniera un plan validado por código,
personajes cuya percepción calcula el código y un narrador que cura su registro sin inventar; los dos
más cercanos, Yu et al. y PLOTTER, reúnen parte de esas piezas. En la evaluación, la literatura revisada desaconseja las escalas
absolutas y los jueces LLM como medida de calidad, y respalda las comparaciones por pares.

El diseño del generador (Capítulo~\ref{chapter:generator}) separa qué pasa, cómo pasa y cómo se
cuenta. Un plan, validado por código como grafo acíclico y con un contrato de promesas anclado a sus
eventos, decide qué pasa; una función de actores que no ven el plan y recuerdan solo lo que
percibieron decide cómo pasa, con un director que sugiere y cierra cada \textit{beat}, y un narrador
que cura el registro sin inventar decide cómo se cuenta. La configuración de referencia mantiene
mecanismos opcionales que pueden activarse para estudiar su efecto sobre las medidas de la prosa y
la preferencia humana.

La metodología de evaluación (Capítulo~\ref{chapter:evaluation}) articula tres evaluaciones en las
que ninguna cifra automática se presenta como juicio de calidad: 24 rasgos núcleo, 18 de ellos
recuentos de un extractor con una cita que el código verifica literalmente y 6 medidas calculadas por
código, tres de ellas sobre lo que extrae el modelo; una evaluación humana por pares a ciegas con
tres preguntas y una línea base de un solo \textit{prompt}, y un juez que intenta reproducir el orden
humano a partir de los rasgos y se valida fuera de muestra, frente al azar y frente a la longitud
sola.

La implementación (Capítulo~\ref{chapter:implementation}) se organizó en cinco paquetes con
dependencias en un solo sentido, de modo que la evaluación no importa el generador y mide del mismo
modo cualquier historia guardada. Cada decisión del código vive en un validador puro con pruebas: el 1
de noviembre de 2026 la suite pasó 907 pruebas, entre ellas las que comprueban que ningún actor ve el
plan y que, con el inventario apagado, el generador queda como estaba. Cada ejecución queda en disco con su
versión y su huella, y las tablas del estudio, del juez y del ranking se regeneran desde esos
artefactos. La cuota diaria
del servicio de modelos fue una restricción práctica que el sistema administra, pero no evita.

La evaluación experimental (Capítulo~\ref{chapter:experiments}) documentó una configuración de
referencia y un conjunto horizontal de 24 ejecuciones: cinco bases simuladas, nueve líneas base,
dos narrativas desde el plan y ocho variantes de narración. La validación registrada del extractor
comprende doce historias, con repetición en seis; sus resultados se presentan por rasgo, con el
número de citas juzgadas y de pares efectivamente comparables. La generación terminada y la
extracción completa se verifican por separado.

El estudio humano, según la actualización del autor, reunió ocho graduados de Ciencia de la
Computación, incluido él mismo, con diferentes frecuencias de lectura. Cada participante comparó
cinco pares de diez textos distintos mediante Telegram y respondió tres preguntas por par, para un
total de 120 respuestas declaradas. Este alcance describe la participación; las preferencias por
criterio y el acuerdo requieren el registro individual, que no se sustituye por los votos del
piloto anterior. Los cinco pares disjuntos tampoco permiten identificar un ranking humano global.

Los contrastes responden a cuestiones diferentes. La matriz conserva escritura directa y simulación
vinculadas por plan; la preferencia entre ambas necesita relacionarse con los pares humanos. El autor
informa de resultados favorables a la memoria propia, cuya magnitud debe documentarse con las
ejecuciones y recuentos correspondientes. Las sesiones conservadas de mecanismos opcionales aportan
una observación compatible con la tercera hipótesis: reducir defectos medidos no aseguró que se
prefiriera el texto modificado. Extender estas observaciones a otras historias y lectores exige una
muestra mayor y la trazabilidad completa de las comparaciones.

Las \textbf{aportaciones} son el generador y la metodología de evaluación. El primero combina
planificación, simulación y narración con comprobaciones explícitas en sus interfaces. La segunda
articula recuentos citados, preferencias humanas por pares y un juez aprendido de esas preferencias.
El rango de acierto comunicado para el juez es de 0,50 a 0,65; sin su desglose y partición no se
interpreta como rendimiento verificado fuera de muestra. El corpus de 122 historias destinado a
ordenar se distingue de los textos que aportan votos al entrenamiento y a la validación.

Las \textbf{limitaciones} se refieren a la escala, la representatividad y la documentación de los
resultados. Los participantes comparten formación y el autor conoce el sistema; las lecturas
adicionales no cuentan como nuevos participantes. La cuota limita las ejecuciones y extracciones,
y los informes históricos no reemplazan los registros del estudio actualizado. Ampliar la muestra,
conectar los pares y conservar las respuestas y predicciones permitirá estimar la incertidumbre.
Usar el juez como función objetivo y buscar configuraciones con un algoritmo genético quedan como
trabajo futuro, una vez demostrada su generalización.

\end{conclusions}

\begin{recomendations}
% Revisión 2026-11-07: resultados actualizados y procedencia en Experiments.tex y fix.md.

La primera recomendación es \textbf{ampliar y diversificar la muestra de lectores}, incorporando
personas con otras formaciones y hábitos de lectura. Conviene conservar las respuestas por lector,
par y criterio, identificar al autor y comparar los resultados con y sin su participación. Alternar
el orden de presentación y la posición A/B permite estudiar los efectos de cansancio y posición.
Las nuevas comparaciones deberían conectar los cinco pares para poder estimar un ranking humano
global, además de preferencias dentro de cada pareja.

Se recomienda \textbf{completar la trazabilidad del estudio}: vincular los diez textos a sus
identificadores de ejecución, recuperar las respuestas individuales y documentar por separado las
lecturas adicionales de los mismos participantes. La comparación de memoria propia y compartida
comunicada por el autor debe acompañarse de los planes utilizados, los recuentos de menciones a
hechos desconocidos y las preferencias. Su replicación, junto con el contraste entre simulación y
escritura directa desde el mismo plan, permitirá estimar la estabilidad de las observaciones.

En el análisis horizontal se debe mantener la distinción entre ejecución completada y extracción
compatible. Es necesario \textbf{reunir las extracciones de las 24 unidades} bajo el mismo protocolo
y conservar las repeticiones que estiman el ruido. Para el extractor, una revisión humana de
historias completas debe registrar tanto las ocurrencias citadas como las omitidas y documentar
quién revisó cada caso. Así podrán contrastarse precisión, cobertura y repetibilidad sin confundir
la ausencia de una anotación con un resultado nulo.

La siguiente prioridad es \textbf{reproducir la validación del juez}: publicar sus aciertos y
denominadores por criterio y partición, conservar las predicciones y compararlo con la longitud
sola sobre los mismos votos. El rango de acierto comunicado debe vincularse con ese informe antes
de atribuirle generalización. Para ordenar el corpus de 122 historias se requieren su lista,
las extracciones compatibles, la identificación de los textos de entrenamiento y las razones de
exclusión; después corresponde analizar cualitativamente las primeras historias.

Como trabajo posterior se propone \textbf{usar el juez como función objetivo}: generar varias
historias con distintos parámetros y seleccionar la que el juez ordene primero, y después buscar
configuraciones mediante un \textbf{algoritmo genético}. Ambos usos requieren un juez que generalice
fuera de muestra y supere al control de longitud. La metodología también puede aplicarse a otros
textos generados, adaptando el catálogo de rasgos y las preguntas, sin sustituir las preferencias
humanas por un juicio de calidad del LLM.

\end{recomendations}

% Margen elástico para las referencias con títulos largos y rangos de páginas (R11).
{\setlength{\emergencystretch}{3em}%
\printbibliography[heading=bibintoc]}


\end{document}